% !TEX program = xelatex
% 电工电子技术 · 第三章第一讲（Chapter Three-1）· 精讲讲义（lecture-deepdive 规范）
\documentclass[11pt]{ctexart}
%=====================================================================
% theme-violet.tex —— 讲义精讲 · 紫花主题
% 用法：\documentclass[11pt]{ctexart}
%       \input{assets/theme-violet}
%       \begin{document} ... \end{document}
% 只用 ctexart 自带能力 + xcolor/enumitem/fancyhdr/pifont/amsmath，
% 不依赖 tikz/tcolorbox，保证 MiKTeX 可直接编译。
%=====================================================================

\usepackage[a4paper,left=2.0cm,right=2.0cm,top=2.1cm,bottom=2.3cm,%
            headsep=12pt,footskip=22pt]{geometry}
\usepackage{amsmath,amssymb}
\usepackage{xcolor}
\usepackage{array}
\usepackage{enumitem}
\usepackage{pifont}
\usepackage{fancyhdr}
\usepackage[hidelinks]{hyperref}

%---------------------------------------------------------------- 配色
\definecolor{violet}{RGB}{107,78,143}      % 主色：紫
\definecolor{violetdeep}{RGB}{63,44,92}    % 深紫：标题/强调
\definecolor{petal}{RGB}{183,159,208}      % 花瓣：装饰线、花饰
\definecolor{bandviolet}{RGB}{247,243,251} % 浅紫底：档案/提示框
\definecolor{bandcream}{RGB}{252,249,242}  % 米白底：易错/注意框
\definecolor{rulepurple}{RGB}{217,207,230} % 分隔线
\definecolor{inkgray}{RGB}{92,96,110}      % 次要文字

%---------------------------------------------------------------- 字体
% 有 Noto CJK 用 Noto，没有就退回 ctex 默认（Windows 下为宋体/黑体）
\IfFontExistsTF{Noto Serif CJK SC}{\setCJKmainfont{Noto Serif CJK SC}}{}
\IfFontExistsTF{Noto Sans CJK SC}{\setCJKsansfont{Noto Sans CJK SC}}{}

%---------------------------------------------------------------- 版式
\setlist{nosep,leftmargin=1.9em}
\linespread{1.28}
\setlength{\parskip}{3.5pt}
\setlength{\parindent}{0pt}
\setlength{\tabcolsep}{5pt}
\renewcommand{\arraystretch}{1.28}

%---------------------------------------------------------------- 花饰
% \ding{95} = 五瓣小花，\ding{96} = 星形花；都来自 pifont，必定存在
\newcommand{\orn}{\textcolor{petal}{\ding{95}}}          % 行内小标
\newcommand{\ornstar}{\textcolor{petal}{\ding{96}}}      % 星形花
\newcommand{\flowerrule}{%                                % 居中花饰分隔线
  \par\vspace{7pt}%
  \noindent\makebox[\linewidth]{\color{petal}%
    \rule{0.20\linewidth}{0.6pt}\hspace{8pt}{\small\ding{95}}\hspace{4pt}%
    \raisebox{1.1pt}{\tiny\ding{95}}\hspace{4pt}{\small\ding{95}}\hspace{8pt}%
    \rule{0.20\linewidth}{0.6pt}}%
  \par\vspace{7pt}}

%---------------------------------------------------------------- 标题与节头
\newcommand{\doctitle}[3]{%   #1 主标题  #2 副标题  #3 元信息（可用 \\ 分行）
  \begin{center}
    {\color{petal}\small \orn\enspace #2\enspace \ornstar}\\[7pt]
    {\LARGE\bfseries\color{violetdeep} #1}\\[10pt]
    {\color{petal}\rule{0.55\linewidth}{1pt}}\\[4pt]
    {\small\color{inkgray} #3}
  \end{center}\vspace{2pt}}

\newcommand{\hdr}[2]{%   #1 章节号（如 §0）  #2 标题
  \par\vspace{12pt}%
  \noindent\colorbox{violetdeep}{\color{white}\bfseries\small\;#1\;}\hspace{7pt}%
  {\color{violetdeep}\large\bfseries #2}\par\nobreak
  \vspace{3pt}{\color{petal}\hrule height 1.2pt}\vspace{6pt}}

\newcommand{\subhdr}[1]{%
  \par\vspace{6pt}%
  \noindent{\color{violet}\bfseries #1}\par\vspace{1pt}}

%---------------------------------------------------------------- 提示框
\newcommand{\callout}[1]{% 浅紫底：讲义档案、读法说明
  \par\medskip\noindent
  \fcolorbox{rulepurple}{bandviolet}{\parbox{\dimexpr\linewidth-2\fboxsep-2\fboxrule}{\vspace{2pt}#1\vspace{2pt}}}%
  \par\medskip}

\newcommand{\warnbox}[1]{% 米白底：易错、边界条件
  \par\medskip\noindent
  \fcolorbox{rulepurple}{bandcream}{\parbox{\dimexpr\linewidth-2\fboxsep-2\fboxrule}{\vspace{2pt}#1\vspace{2pt}}}%
  \par\medskip}

\newcommand{\gapitem}[2]{% #1 缺口类型（紫） #2 内容
  \par\vspace{2pt}\noindent
  \hangindent=1.4em\hangafter=1
  {\color{petal}\small\ding{96}}\enspace{\bfseries\color{violetdeep}#1}\enspace #2\par}

% 前提旁注：只放在「需要理解的那一步」后面，给陈述 + 本题代入，不重证。
\newcommand{\aside}[1]{%
  \par\vspace{1pt}\noindent
  \hangindent=1.6em\hangafter=1
  {\small\color{inkgray}%
    \hspace*{0.6em}{\color{petal}$\triangleright$}\hspace{2pt}#1}%
  \par\vspace{2pt}}

%---------------------------------------------------------------- 表格
% 本机 MiKTeX 的 colortbl 与 array 版本不匹配（\rowcolor 会报错），
% 故表格统一用「上粗紫线 + 表头 + 中细线」的三线式，不依赖 colortbl。
\newcommand{\tabtop}{\par\vspace{3pt}{\color{violet}\hrule height 1.2pt}\vspace{1pt}}
\newcommand{\tabmid}{\par\vspace{1pt}{\color{petal}\hrule height 0.5pt}\vspace{1pt}}
\newcommand{\tabbot}{\par\vspace{3pt}{\color{rulepurple}\hrule height 0.5pt}\vspace{8pt}}
\newcommand{\thead}[1]{{\bfseries\color{violetdeep}#1}}

%---------------------------------------------------------------- 页眉页脚
\pagestyle{fancy}
\fancyhf{}
\renewcommand{\headrulewidth}{0.4pt}
\renewcommand{\headrule}{\hbox to\headwidth{\color{rulepurple}\leaders\hrule height \headrulewidth\hfill}}
\fancyhead[L]{\small\color{inkgray}\songti \doctag}
\fancyhead[R]{\small\color{inkgray}\songti \coursetag}
\fancyfoot[C]{\small\color{inkgray}\orn\enspace{\color{violet}\thepage}\enspace\orn}

% 每篇讲义在导言区设这两个值
\newcommand{\doctag}{}
\newcommand{\coursetag}{}
\newcommand{\settags}[2]{\renewcommand{\doctag}{#1}\renewcommand{\coursetag}{#2}}

%---------------------------------------------------------------- 数学记号
\newcommand{\vect}[1]{\boldsymbol{#1}}
\newcommand{\dd}{\,\mathrm{d}}
\newcommand{\ddx}{\frac{\mathrm{d}}{\mathrm{d}x}}
\addtolength{\jot}{2pt}



\settags{C3-1 · 半导体基础与二极管 · 精讲}{电工电子技术}

\begin{document}

\doctitle{半导体基础与二极管：PN 结如何变成一只单向开关}
{从一根断键、一个势垒，到三种等效模型——3.1--3.2 全册精讲}
{依据：《Chapter Three-1-2026 中白上传》PDF（p1--p69，共 69 页；本讲义页码一律指本册 PDF 物理页）\ \textperiodcentered\ 2026-10-23}

%=============================================================
\hdr{§0}{课件坐标、记号与三处拼接}

\subhdr{课件讲了哪几块，各在哪几页}

这册课件是\textbf{一条因果链}：从「一块硅里有多少能动的电荷」出发（3.1），
把两种掺杂的硅拼在一起得到 PN 结，让结自己长出一个内建电场（3.1.4），
再加上外电压看它导通还是截止（单向导电），最后把整块结封装成两端的器件——二极管，
用一条伏安特性曲线和三种等效模型去用它（3.2）。
链上每一步课件都画了图，但\textbf{每一步的「为什么」几乎都留白}：全册 69 页，公式只有零星的
$U_D=0.7\,\mathrm{V}$、$U_F=U_B+U_D$ 几处，连二极管电流的指数式都没有写出来。

\par\vspace{2pt}
\noindent\begin{tabular}{@{}>{\raggedright\arraybackslash}p{1.4cm}>{\raggedright\arraybackslash}p{3.0cm}>{\raggedright\arraybackslash}p{6.9cm}>{\raggedright\arraybackslash}p{4.2cm}@{}}
\tabtop
\thead{课件页} & \thead{小节} & \thead{内容} & \thead{本讲义位置}
\\
\tabmid
p1 & 封面 & Chapter Three DC regulated power supply & §0
\\
p2 & 3.1.1 & 导体/半导体/绝缘体三类材料对比 & §1.1
\\
p3--p8 & 3.1.2 I、II & 共价键、8 电子、两种载流子、$n=p$ & §1.2、§1.3
\\
p9--p11 & 3.1.2 II & 载流子漂移方向、本征导电弱、掺杂出口 & §1.3、§1.4
\\
p12--p17 & 3.1.3 I、II & 掺杂定义、磷（+5）、多余电子、多子 & §2.1
\\
p18--p19 & 3.1.3 II & $n\gg p$、多数/少数载流子 & §2.1、§2.3
\\
p20--p24 & 3.1.3 III & 硼（+3）、空穴、多子、导电能力提高 & §2.2、§2.4
\\
p25--p26 & 3.1.3 IV & 两幅晶格图的判别题（只问不答） & §2.5
\\
p27--p32 & 3.1.4 I & 两块材料、拼合、扩散/漂移两过程 & §3.1
\\
p33--p36 & 3.1.4 I & 扩散、固定离子、内建电场、动态平衡 & §3.2--§3.5
\\
p37--p40 & 3.1.4 II & 正偏、耗尽层变薄、约 $0.7\,\mathrm{V}$ 导通 & §4.1
\\
p41--p43 & 3.1.4 II、III & 反偏、耗尽层变厚、ON/OFF 总结 & §4.2、§4.3
\\
p44--p48 & 3.2、3.2.1 I & 二极管结构、符号、封装、分类 & §5.1
\\
p49--p50 & 3.2.1 II-1 & 伏安特性曲线、五个定义性问句 & §5.2--§5.4
\\
p51--p52 & 3.2.1 II-2、3 & 理想模型、恒压降模型 & §6
\\
p53--p55 & 3.2.1 II & 温度两条曲线、Si/Ge 两条曲线（提问） & §7.1、§7.2
\\
p56 & 3.2.1 III & 四个主要参数 & §7.3
\\
p57--p61 & 3.2.4 Clamping & 钳位电路五个追问 & §8
\\
p62--p69 & 3.2.4 E1--E3 & 二极管逻辑、真值表逐页填 & §9
\tabbot
\end{tabular}

\subhdr{本册与前后册的接口}

这本 C3-1 停在「二极管 + 三种模型」，但它的标题是 「DC regulated power supply」（p1）。
后面两册的分工是：
（1）C3-2：把二极管接成整流电路（半波、桥式），把交流变成脉动直流；
（2）C3-3：加滤波电容与稳压环节，把脉动直流变成稳定直流。
所以本册的每一个结论在那里都要用一次：$0.7\,\mathrm{V}$ 决定整流管的导通角，
$U_{BR}$ 决定选管的耐压，$I_S$ 与温度关系决定滤波后的纹波与热设计。
读本册时把「这是后面哪一步的前提」记住，比背公式有用。

\subhdr{记号：先统一，后文不再回头}

\par\vspace{1pt}{\small
\noindent\begin{tabular}{@{}>{\raggedright\arraybackslash}p{2.6cm}>{\raggedright\arraybackslash}p{11.4cm}@{}}
\tabtop
\thead{记号} & \thead{约定}
\\
\tabmid
$U_D$、$I_D$ & 二极管\textbf{两端}的直流电压与电流；参考方向一律取\textbf{阳极}$\to$\textbf{阴极}，即正向为 $U_D>0$、$I_D>0$
\\
$U_F$、$I_F$ & 正向压降、额定正向平均电流（p56 的器件参数）；$U_F$ 是「导通时管压降」的通用记号
\\
$U_R$、$U_{BR}$、$I_{Rm}$ & 最高反向工作电压、反向击穿电压、最大反向电流（p56）
\\
$U_T$ & 热电压 $U_T=kT/q$，300 K 时 $25.9\,\mathrm{mV}$（工程上取 $26\,\mathrm{mV}$）
\\
$n_i$ & 本征载流子浓度；$N_A$、$N_D$ 受主、施主浓度
\\
$E_g$ & 禁带宽度（eV）；$q$ 电子电荷量，$k$ 玻尔兹曼常数
\\
$r_d$ & 二极管的\textbf{动态}（微分）电阻 $\mathrm{d}U_D/\mathrm{d}I_D$
\\
$V_{bi}$ & 内建电势（势垒高度除以 $q$），单位 V
\tabbot
\end{tabular}}

\aside{大写 $U,I$ 一律指直流量；本册不涉及交流小信号，出现的小写只在讨论动态电阻时使用。}

\subhdr{三处拼接（课件自己的缝）}

\gapitem{拼接缺口}{p11 用一行箭头把「Intrinsic semiconductor $\to$ Extrinsic semiconductor」带过，
却没有回答它自己提的问题（p11 的 (9) 「How to increase the conductivity?」）。
从「载流子少」到「掺一点杂质就够」中间缺一个\textbf{数量级}：本征 $n_i\approx1.5\times10^{10}\,\mathrm{cm^{-3}}$，
掺杂 $10^{16}\,\mathrm{cm^{-3}}$，差六个数量级。本讲义 §2.3--§2.4 把这条缝补成算式。}

\gapitem{拼接缺口}{p43 以 ON/OFF 收束「单向导电」，p49 突然出现一条平滑的伏安曲线，
两页之间没有任何过渡。缺的是「开关」与「曲线」的桥：曲线正是把 ON/OFF 的突变
展开成\textbf{指数过渡}的结果。§5.2 用肖克利方程把桥搭上。}

\gapitem{形式改写}{p52 把硅管的导通压降写成常数 $U_D=0.7\,\mathrm{V}$，
而 p55 的锗管曲线拐点读出来约 $0.2\,\mathrm{V}$，与 p52 写的 $0.3\,\mathrm{V}$ 对不上。
两个数都不是错，差别在于「读图的位置」：$0.7/0.3\,\mathrm{V}$ 是\textbf{mA 级电流}下的工程约定，
曲线拐点则是电流刚开始明显上升的位置。§6.2 统一口径，§7.2 点名这处不一致。}

\subhdr{三条主线：这册只在回答三个问题}

（1）\textbf{有多少能动的电荷}（3.1.1--3.1.3）：本征只有 $n_i$，掺杂把它抬到 $10^{16}$ 量级——
这是「材料」问题，答案是一串浓度；
（2）\textbf{两种电荷拼在一起会长出什么}（3.1.4）：扩散留下固定离子，固定离子产生内建电场，
内建电场又拦住扩散，最后停在动态平衡——这是「机制」问题，答案是 $V_{bi}$、$W$、$E_{\max}$；
（3）\textbf{两端看进去像什么元件}（3.2）：把整块结封装成二极管，用一条曲线、三个模型、四个参数描述——
这是「器件」问题，答案是 $I=I_S(e^{U/(nU_T)}-1)$ 与 $0.7\,\mathrm{V}$。

三个问题对应三种问法：问「N 型多子是什么」考第一层，问「为什么加 $0.7\,\mathrm{V}$ 才导通」考第二层，
问「用理想模型还是恒压降模型算」考第三层。三层的答案不能互相代替。

\callout{\textbf{本讲义怎么用}：\\
读第一遍只看 §0 表、§6 的模型对照表、§8 的钳位四步、§9 的取最低/取最高口诀——这四处直接决定题能不能做对；\\
读第二遍看 §3.3、§4.1、§5.2——这三处解释 $0.7\,\mathrm{V}$ 从哪来，是「换一只管子、换一个温度还能不能算对」的关键；\\
§10 缺口清单逐条对应课件某页某句，卡住时去那里找。}

%=============================================================
\hdr{§1}{半导体为什么「半」导电：禁带、共价键与两种载流子}

\subhdr{1.1 三类材料的差别只有一个来源：禁带宽度（p2）}

课件 p2 用一张分类图和四行英文把三类材料并排：
导体「Low resistance. Always conduct electric current.」；
半导体「Conductivity in-between. Block or allow the current flow, providing total control over it.」；
绝缘体「Very high electrical resistance. Do not allow the flow of current.」。
它没说差别从哪来。差别只有一个来源：\textbf{价带顶与导带底之间的禁带宽度 $E_g$}。

\par\vspace{2pt}
\noindent\begin{tabular}{@{}>{\raggedright\arraybackslash}p{2.6cm}>{\raggedright\arraybackslash}p{2.9cm}>{\raggedright\arraybackslash}p{5.3cm}>{\raggedright\arraybackslash}p{4.2cm}@{}}
\tabtop
\thead{材料} & \thead{$E_g$（300 K）} & \thead{载流子来源} & \thead{室温电导率}
\\
\tabmid
铜 Cu & $0$（价带与导带重叠） & 自由电子气，密度 $\sim10^{23}\,\mathrm{cm^{-3}}$ & $5.8\times10^{5}\,\mathrm{S/cm}$
\\
硅 Si & $1.12\,\mathrm{eV}$ & 热激发断键，密度 $\sim10^{10}\,\mathrm{cm^{-3}}$ & $4.4\times10^{-6}\,\mathrm{S/cm}$
\\
锗 Ge & $0.67\,\mathrm{eV}$ & 同上，但禁带窄、断键多 & $\sim2\times10^{-2}\,\mathrm{S/cm}$
\\
二氧化硅 & $\sim9\,\mathrm{eV}$ & 常温下几乎为零 & $<10^{-16}\,\mathrm{S/cm}$
\tabbot
\end{tabular}

\textbf{为什么禁带决定一切}：电子要参与导电，必须从价带跳到导带。跳过去所需的能量是 $E_g$，
而常温提供的热能量级只有 $kT\approx25.9\,\mathrm{meV}$。对硅，
\[
\frac{E_g}{2kT}=\frac{1.12}{2\times0.0259}\approx21.7,
\qquad
e^{-21.7}\approx3.7\times10^{-10}.
\]
也就是说，只有约 $10^{-10}$ 这个量级的电子能靠热激发跨过禁带。
这个极小的因子，就是「半导体」三个字的全部含义：\textbf{能导电，但载流子少得可怜}。
[教材补充：本小节与下一小节的能带、态密度说法不在课件里，补它是因为 p2 的四行英文需要一个定量解释。]

\subhdr{1.2 本征半导体：八个电子、一根键、成对产生的载流子（p3--p8）}

p3 的 (2) 问 「What are semiconductor materials?」，p3 底部答 「Silicon and germanium are the
most widely used intrinsic semiconductors.」。p4 的 (3) 问 「Why there are eight electrons outside?」，
答案是两句：每个硅原子最外层有 4 个价电子（「four electrons (valence electrons) in its outermost shell」），
每个原子与四个邻居各共用一对（「Each atom shares its four valence electrons with its four neighbors」）。
$4+4=8$，凑满最外层八电子结构（p5：「An atom is most stable when it has eight electrons in its
outermost energy level」）。

一根共价键被热激发打断时，会发生两件事，而且\textbf{必然同时发生}：
（1）一个价电子挣脱束缚变成自由电子；（2）它原来的位置留下一个空位，即空穴。
p7 把这两半分别命名为 「Free electron」 与 「hole」，p8 的 (6) 问两者数目是否相同，
答 「The number of holes and free electrons in an intrinsic semiconductor are equal.」——
\textbf{相等不是巧合，是「成对产生」的直接结果}：断一根键只产生一个电子加一个空穴，
复合时也是一个电子填一个空穴。若记本征载流子浓度为 $n_i$，则恒有
\[
n=p=n_i .
\]

\begin{center}
\begin{minipage}{0.94\linewidth}
\subhdr{质量作用定律（原形式）}
设半导体处于热平衡、非简并，导带有效态密度为 $N_C$、价带有效态密度为 $N_V$，禁带宽度 $E_g$。
则电子浓度 $n$ 与空穴浓度 $p$ 的\textbf{乘积}只由材料与温度决定：
\[
np=N_C N_V\,e^{-E_g/(kT)}=n_i^2 ,
\]
与掺杂浓度无关。其中 $N_C,N_V\propto T^{3/2}$，故
\[
n_i=\sqrt{N_C N_V}\;e^{-E_g/(2kT)}\ \propto\ T^{3/2}e^{-E_g/(2kT)} .
\]
\aside{课件没有写出这条定律，p8 只给了本征时的特例 $n=p$。本讲义把它写成一般形式，
因为 §2 的每一次「多子/少子」计算都要用它；$n_i$ 的数值：Si $1.5\times10^{10}\,\mathrm{cm^{-3}}$、
Ge $2.4\times10^{13}\,\mathrm{cm^{-3}}$（300 K）。}
\end{minipage}
\end{center}

\textbf{把数量级摆出来}：硅原子密度约 $5\times10^{22}\,\mathrm{cm^{-3}}$，本征载流子 $1.5\times10^{10}\,\mathrm{cm^{-3}}$，
平均每
\[
\frac{5\times10^{22}}{1.5\times10^{10}}\approx3.3\times10^{12}
\]
个硅原子才有一根键断开。p10 说 「The carriers of intrinsic semiconductor are few in number.
So the conductive ability is weak.」——「few」是 $3\times10^{12}$ 分之一，这就是不掺杂的硅几乎不导电的原因。

\subhdr{1.3 两种载流子怎么动、电流往哪边流（p7--p9）}

p9 的 (7) 问 「Do free electrons and holes move in the same direction when there is a potential?」，
答句是 「When a potential difference is applied electrons move towards the higher potential and
holes move towards the lower potential.」。这句话里藏着两个容易混的点。

\textbf{第一，空穴为什么会「动」。}空穴本身不是粒子，它只是共价键上的一个缺位。
当相邻键上的价电子跳过来填补这个缺位时，缺位就转移到了旁边——宏观上看起来像是一个带正电的
「空穴」在朝相反方向移动。所以空穴的移动\textbf{本质上是价电子的反向接力}，
它的电荷取 $+q$、有效质量为正，可以当成一种准粒子来算。

\textbf{第二，电流方向由正电荷定义。}电子带负电，向高电位运动，等效电流方向与运动方向相反；
空穴带正电，向低电位运动，等效电流方向与运动方向相同。两种载流子运动方向相反，
但它们贡献的电流方向\textbf{相同}——都是从高电位流向低电位。

\aside{p9 的图 2 把这件事画成两排符号：左排电子「$-$」向右移、下方电流箭头向左；
右排空穴「$+$」向右移、下方电流箭头向右。代入本册记号：电流正方向取阳极$\to$阴极，
与空穴运动同向、与电子运动反向。}

两种载流子对电流的贡献可以直接相加：
\[
I=qA\left(n v_n+p v_p\right),
\]
其中 $v_n,v_p$ 是漂移速度。电场不太强时 $v=\mu E$，于是
\[
I=qA\left(n\mu_n+p\mu_p\right)E
\quad\Longrightarrow\quad
\sigma=\frac{I}{AE}=q\left(n\mu_n+p\mu_p\right).
\]

\subhdr{1.4 本征导电有多弱：把 $\sigma$ 算出来（p10）}

取硅在 300 K 的迁移率 $\mu_n=1350\,\mathrm{cm^2/(V\cdot s)}$、$\mu_p=480\,\mathrm{cm^2/(V\cdot s)}$，
本征时 $n=p=n_i=1.5\times10^{10}\,\mathrm{cm^{-3}}$，代入上式
\[
\sigma_i=q\,n_i\left(\mu_n+\mu_p\right)
=1.6\times10^{-19}\times1.5\times10^{10}\times1830
\approx4.4\times10^{-6}\,\mathrm{S/cm}.
\]
对应的电阻率 $\rho_i=1/\sigma_i\approx2.3\times10^{5}\,\Omega\cdot\mathrm{cm}$，
比铜（$1.7\times10^{-6}\,\Omega\cdot\mathrm{cm}$）大了约 $10^{11}$ 倍。
p10 的 「conductive ability is weak」 到这里有了数字：\textbf{弱十一个数量级}。

p11 的 (9) 问 「How to increase the conductivity?」，课件只给了一个箭头：
Intrinsic $\to$ Extrinsic。为什么不掺杂不行、掺了又能好多少，下一节算。

%=============================================================
\subhdr{1.5 本征浓度随温度跑得有多快}

$n_i\propto T^{3/2}e^{-E_g/(2kT)}$ 里指数项压倒一切。把硅的数值列出来：

\par\vspace{2pt}
\noindent\begin{tabular}{@{}>{\raggedright\arraybackslash}p{3.4cm}>{\raggedright\arraybackslash}p{4.2cm}>{\raggedright\arraybackslash}p{3.4cm}>{\raggedright\arraybackslash}p{3.8cm}@{}}
\tabtop
\thead{温度} & \thead{$n_i$（$\mathrm{cm^{-3}}$）} & \thead{与 300 K 之比} & \thead{含义}
\\
\tabmid
250 K（$-23\,{}^\circ\mathrm{C}$） & $1.5\times10^{8}$ & $10^{-2}$ & 载流子少，杂质也未全电离
\\
300 K（$27\,{}^\circ\mathrm{C}$） & $1.5\times10^{10}$ & $1$ & 室温基准
\\
350 K（$77\,{}^\circ\mathrm{C}$） & $4.2\times10^{11}$ & $28$ & 漏电开始明显
\\
400 K（$127\,{}^\circ\mathrm{C}$） & $5.2\times10^{12}$ & $347$ & 硅器件的结温上限附近
\\
500 K（$227\,{}^\circ\mathrm{C}$） & $1.9\times10^{14}$ & $1.2\times10^{4}$ & 接近轻掺杂浓度
\\
600 K（$327\,{}^\circ\mathrm{C}$） & $2.1\times10^{15}$ & $1.4\times10^{5}$ & 与 $10^{15}$ 掺杂可比
\tabbot
\end{tabular}

\textbf{两个推论}：
（1）硅器件的结温上限（常见 $125$--$150\,{}^\circ\mathrm{C}$）\textbf{不是因为}这块硅快变成本征了
（那要 $700\,\mathrm{K}$，§2.3），而是因为 $n_i$ 在 $150\,{}^\circ\mathrm{C}$ 已涨到 $1.4\times10^{13}$，
反向漏电与功耗随之上升，容易热失控；
（2）任何「$n\approx N_D$」的近似都默认 $N_D\gg n_i(T)$。温度一高，这个近似先失效，
而不是先「烧管子」。

\subhdr{1.6 为什么偏偏是硅和锗：一张材料对照表}

p3 说 「Silicon and germanium are the most widely used intrinsic semiconductors.」，
没说为什么。三种常见半导体的关键量并排：

\par\vspace{2pt}
\noindent\begin{tabular}{@{}>{\raggedright\arraybackslash}p{2.2cm}>{\raggedright\arraybackslash}p{2.6cm}>{\raggedright\arraybackslash}p{2.6cm}>{\raggedright\arraybackslash}p{3.0cm}>{\raggedright\arraybackslash}p{3.2cm}@{}}
\tabtop
\thead{材料} & \thead{$E_g$/eV} & \thead{$\mu_n$} & \thead{$\mu_p$} & \thead{器件温度上限}
\\
\tabmid
Si & $1.12$ & $1350$ & $480$ & $\sim150\,{}^\circ\mathrm{C}$
\\
Ge & $0.67$ & $3900$ & $1900$ & $\sim75$--$90\,{}^\circ\mathrm{C}$
\\
GaAs & $1.42$ & $8500$ & $400$ & $\sim200\,{}^\circ\mathrm{C}$
\tabbot
\end{tabular}
（迁移率单位 $\mathrm{cm^2/(V\cdot s)}$，300 K。）

读表：
（1）\textbf{锗的迁移率高、拐点低}，所以同样电流下压降小、更适合低压；但 $E_g$ 小，
$n_i$ 大两个数量级，反向漏电大、温度上限低——这就是锗被硅取代的原因；
（2）\textbf{硅的 $E_g$ 居中}，$150\,{}^\circ\mathrm{C}$ 仍可用，且能长出优质 $\mathrm{SiO_2}$ 绝缘层，
工艺成熟——三项加起来才是硅的主流地位；
（3）\textbf{砷化镓迁移率最高}，适合高频，但 $E_g$ 大、拐点约 $1.4\,\mathrm{V}$，且成本高。
[教材补充：材料对照与温度上限不在课件中；p45 的目录里有 Photodiode 与 LED，
本册没有展开，它们在后续册里讲。]

\hdr{§2}{杂质半导体：掺杂改的是「谁多谁少」}

\subhdr{2.1 施主与 N 型：第五个价电子（p12--p19）}

p12 的 (2) 给的定义是 「Doping: the process by which an impurity is added to a semiconductor」。
p13 按杂质类型分成 N 型与 P 型。p14 的 (4) 给出 N 型的来源：
「By adding some phosphorus(P) with a valence of +5.」——把硅晶格里的一个硅原子换成磷原子。

磷有 5 个价电子，而晶格只要求它出 4 个参与共价键。第 5 个电子没有键位可占，
只受磷原子核的库仑吸引，束缚能只有约 $45\,\mathrm{meV}$（[教材补充：硅中磷的电离能]），
远小于 $kT=25.9\,\mathrm{meV}$ 的几倍，所以\textbf{室温下几乎全部电离}：
\[
\mathrm{P}\ \longrightarrow\ \mathrm{P}^{+}+e^{-}.
\]
电离后，磷原子变成固定在晶格上的正离子（$+q$，不能移动），多出来的电子成为自由电子。
这就是 p16 的 「surplus electron」 与 「Can be a free electron at normal temperature」。

设施主浓度 $N_D$（单位 $\mathrm{cm^{-3}}$），并假设全部电离。电中性要求
\[
n=N_D+p .
\]
在 N 型里 $N_D\gg n_i$，所以 $n\approx N_D$，而空穴靠热激发产生，由质量作用定律
\[
p=\frac{n_i^2}{n}\approx\frac{n_i^2}{N_D}.
\]

\par\vspace{2pt}
\noindent\begin{tabular}{@{}>{\raggedright\arraybackslash}p{3.0cm}>{\raggedright\arraybackslash}p{4.0cm}>{\raggedright\arraybackslash}p{4.0cm}>{\raggedright\arraybackslash}p{3.6cm}@{}}
\tabtop
\thead{施主浓度 $N_D$} & \thead{$n\approx N_D$} & \thead{$p\approx n_i^2/N_D$} & \thead{$n/p$}
\\
\tabmid
$10^{14}\,\mathrm{cm^{-3}}$ & $10^{14}$ & $2.3\times10^{6}$ & $4.4\times10^{7}$
\\
$10^{16}\,\mathrm{cm^{-3}}$ & $10^{16}$ & $2.3\times10^{4}$ & $4.4\times10^{11}$
\\
$10^{18}\,\mathrm{cm^{-3}}$ & $10^{18}$ & $2.3\times10^{2}$ & $4.4\times10^{15}$
\tabbot
\end{tabular}

p18 的 (6) 答 「The number of free electrons >> The number of holes」，p19 的 (7) 把自由电子称为
多数载流子（majority carriers）、空穴称为少数载流子（minority carriers）。
「多数」与「少数」不是两个并列的名称，而是\textbf{同一句质量作用定律的两个输出}：
一个被掺杂抬上去，另一个就被按下去，乘积恒为 $n_i^2$。

\subhdr{2.2 受主与 P 型：少一个电子就多一个空穴（p20--p24）}

p20 的 (8) 给出 P 型：「By adding some Boron(B) with a valence of +3.」。
硼只有 3 个价电子，却要参与 4 条共价键，于是有一条键缺一个电子——
这个缺位就是空穴。硼原子接受一个电子后变成固定在晶格上的负离子 $\mathrm{B}^{-}$。
p20 图上标的就是 $\mathrm{B}^{-}$，这是一个容易看漏的细节：\textbf{掺进来的原子本身带电，
但它不能移动}。

同理，设受主浓度 $N_A$，电中性给出
\[
p=N_A+n,\qquad p\approx N_A,\qquad n\approx\frac{n_i^2}{N_A}.
\]
p23 的 (10) 答 「Holes are majority carriers. Free electrons are minority carriers.」，
与 N 型正好互换。

\subhdr{2.3 「远大于」的条件，以及它什么时候失效（p18）}

p18 只写了 $n\gg p$，没有写这个「远大于」依赖什么。依赖的是
\[
N_D\gg n_i .
\]
掺杂浓度比本征浓度高几个数量级，这个不等式才成立。它两个方向都会失效：

（1）\textbf{高温失效}：$n_i$ 随温度指数增长，涨到与 $N_D$ 可比时，掺杂提供的那点载流子被
本征激发淹没，N 型「退回」本征。以 $N_D=10^{16}\,\mathrm{cm^{-3}}$ 为例，数值计算给出
$n_i$ 在 600 K 约 $2.1\times10^{15}$、700 K 约 $1.3\times10^{16}\,\mathrm{cm^{-3}}$，
即\textbf{约 700 K（$427\,{}^\circ\mathrm{C}$）以上}，这块硅就不再是「N 型」了。

（2）\textbf{低温失效}：温度太低时施主电离能（$45\,\mathrm{meV}$）都供不起，
杂质原子不电离，载流子反而减少，这叫冻结（freeze-out）。[教材补充：这两条边界课件未提，
但 p18 的「$\gg$」是一个只在室温附近成立的结论，必须标出适用范围。]

\subhdr{2.4 掺杂把电导率抬高了多少（p24）}

p24 的 (11) 答 「Conductivity is increased.」。抬高多少？N 型中 $n\approx N_D=10^{16}\,\mathrm{cm^{-3}}$，
空穴项可以忽略，
\[
\sigma_n=qN_D\mu_n=1.6\times10^{-19}\times10^{16}\times1350
\approx2.16\,\mathrm{S/cm},
\]
与本征的 $4.4\times10^{-6}\,\mathrm{S/cm}$ 相比，
\[
\frac{\sigma_n}{\sigma_i}\approx4.9\times10^{5}.
\]
电阻率从 $2.3\times10^{5}\,\Omega\cdot\mathrm{cm}$ 降到 $0.46\,\Omega\cdot\mathrm{cm}$。
\textbf{掺杂浓度提高约六个数量级，电导率也提高约六个数量级}——这就是 p11 那个箭头的真正分量。

\subhdr{2.5 p25--p26 判别题的读法}

p25 并排放两幅晶格图，问 「Which one is N type? Which one is P type?」，p26 才给出标注。
判别只需看三层：
（1）圆内符号：$+$ 是施主电离后的正离子（N 型），$-$ 是受主电离后的负离子（P 型）；
（2）圆间小点的形状：实心点是电子，空心圈是空穴；
（3）数量：哪种载流子多，就是哪一型。
p26 的答案与这个顺序一致：左图圆内 $+$、实心点多 $\Rightarrow$ N 型；右图圆内 $-$、空心圈多 $\Rightarrow$ P 型。

三个变式练手：
（1）若把 p25 左图的圆内符号擦掉、只留载流子，则\textbf{实心点多}的一幅是 N 型——因为 N 型多子是电子；
（2）若圆内符号保留、载流子形状对调（实心点变空心圈），则结论以\textbf{圆内符号}为准，
因为固定离子由掺杂元素决定，载流子的画法可能只是排布示意；
（3）若两幅图都不给符号，只给两种浓度，则\textbf{浓度高的那种载流子是多数载流子}，
它所在的材料就是该型。三条都指向同一件事：\textbf{多子是谁，材料就是哪型}。

\subhdr{2.6 掺杂浓度能取多大：从 $10^{14}$ 到 $10^{20}$}

工艺上掺杂浓度大致在 $10^{14}$--$10^{20}\,\mathrm{cm^{-3}}$ 之间，跨六个数量级：
\begin{itemize}
\item $10^{14}$--$10^{15}$：轻掺杂，做高压管的漂移区或高阻衬底；
\item $10^{16}$--$10^{17}$：常用范围，做普通二极管与三极管的基区、发射区；
\item $10^{18}$--$10^{20}$：重掺杂，做欧姆接触，或做稳压管（要薄耗尽层，§4.5）。
\end{itemize}

掺杂浓度直接决定三件事：
（1）多子浓度 $n\approx N_D$ 与少子浓度 $p\approx n_i^2/N_D$（§2.1）；
（2）内建电势 $V_{bi}=U_T\ln(N_A N_D/n_i^2)$：两边都加重，$V_{bi}$ 只按对数慢慢涨；
（3）耗尽层宽度 $W\propto1/\sqrt{N}$：掺杂越重，$W$ 越薄，击穿电压越低。
所以「耐高压」与「低导通电阻」互相矛盾——这正是功率器件设计的核心取舍。
[教材补充：掺杂范围与取舍不在课件中，补它是为了解释为什么 p46 的整流管与 p52 的 $0.7\,\mathrm{V}$ 都默认「中等掺杂」。]

%=============================================================
\subhdr{2.7 爱因斯坦关系：迁移率与扩散系数是同一个数}

§1.3 用迁移率 $\mu$ 写漂移，§3.2 用扩散系数 $D$ 写扩散，两者看似独立，其实由同一个物理量联系。
平衡时净电流为零，即漂移与扩散严格抵消：
\[
q n\mu_n E+qD_n\frac{\mathrm{d}n}{\mathrm{d}x}=0 .
\]
热平衡下 $n\propto e^{q\varphi/(kT)}$，故
\[
\frac{\mathrm{d}n}{\mathrm{d}x}=n\,\frac{q}{kT}\frac{\mathrm{d}\varphi}{\mathrm{d}x}
=-n\,\frac{qE}{kT},
\]
代回即得
\[
D_n=\frac{kT}{q}\mu_n=\mu_nU_T,\qquad D_p=\mu_pU_T .
\]
这条爱因斯坦关系说明：载流子的「好动程度」只有一个参数 $\mu$，扩散与漂移是同一件事的两种看法。
硅的数值：$D_n=1350\times0.0259\approx34.9\,\mathrm{cm^2/s}$、
$D_p=480\times0.0259\approx12.4\,\mathrm{cm^2/s}$。
\aside{本式在 §3.3 求 $V_{bi}$ 时直接用了一次：把 $D_n$ 换成 $\mu_nU_T$，
才能把方程整理成 $\mathrm{d}\varphi=U_T\mathrm{d}n/n$。}

\subhdr{2.8 掺杂对 $V_{bi}$ 与 $W$ 的敏感性：一个对数、一个平方根}

把 $N_A=N_D=N$ 代入 §3.3 的 $V_{bi}$ 与 §3.4 的 $W$（$U=0$），逐个算出来：

\par\vspace{2pt}
\noindent\begin{tabular}{@{}>{\raggedright\arraybackslash}p{3.2cm}>{\raggedright\arraybackslash}p{3.6cm}>{\raggedright\arraybackslash}p{3.6cm}>{\raggedright\arraybackslash}p{2.6cm}@{}}
\tabtop
\thead{$N$（$\mathrm{cm^{-3}}$）} & \thead{$V_{bi}$} & \thead{$W$} & \thead{$W$ 相对上一行}
\\
\tabmid
$10^{14}$ & $0.455\,\mathrm{V}$ & $3.43\,\mu\mathrm{m}$ & ——
\\
$10^{15}$ & $0.574\,\mathrm{V}$ & $1.22\,\mu\mathrm{m}$ & $0.36$
\\
$10^{16}$ & $0.693\,\mathrm{V}$ & $0.42\,\mu\mathrm{m}$ & $0.35$
\\
$10^{17}$ & $0.812\,\mathrm{V}$ & $0.145\,\mu\mathrm{m}$ & $0.34$
\\
$10^{18}$ & $0.932\,\mathrm{V}$ & $0.049\,\mu\mathrm{m}$ & $0.34$
\tabbot
\end{tabular}

\textbf{读表两句话}：
（1）掺杂每加 10 倍，$V_{bi}$ 只涨约 $0.12\,\mathrm{V}$（对数关系），所以「靠掺杂调 $0.7\,\mathrm{V}$」几乎调不动；
（2）掺杂每加 10 倍，$W$ 缩到约三分之一（平方根关系），所以耐压与结电容对掺杂非常敏感。
这解释了 §4.5 的选择：要击穿电压高就轻掺杂（$W$ 厚），要结电容小就重掺杂（$W$ 薄），两者不可兼得。

\hdr{§3}{PN 结：扩散、漂移与动态平衡}

\subhdr{3.1 从两块材料到一条界面（p27--p32）}

p27--p30 反复强调 「Two separate pieces of silicon, one is N-type, the other is P-type.」，
并用三维方块画了两块尚未接触的材料。p31 只加一句 「Now the n-type and the p-type materials
are brought together」，p32 随即把载流子的运动方式分成两类：
「Two basic processes for carrier transport in semiconductor: Diffusion motion and drift motion.」。

p28 有一句容易被当成万能前提的话：「Ignore minority carriers」。它\textbf{只在这一步成立}：
算结区空间电荷、算内建电场时，少子数量少到可以忽略；但 §4.2 算反向电流时，
反向电流恰恰全部由少子漂移提供。把这句话当成「全局忽略少子」，后面就解释不了反向饱和电流。
这是本册第一处适用范围陷阱，§10 记一条。

\subhdr{3.2 扩散：浓度差把多子推过界面（p33--p34）}

p33 的定义是 「Diffusion motion: The movement of the majority carriers because of the
difference of concentration.」。N 型一侧电子浓度 $10^{16}$ 量级，P 型一侧只有 $10^{4}$ 量级，
浓度差六个数量级以上，电子必然向 P 型一侧扩散；空穴反过来向 N 型一侧扩散。
扩散电流与浓度梯度成正比（一维）：
\[
J_n=qD_n\frac{\mathrm{d}n}{\mathrm{d}x},\qquad J_p=-qD_p\frac{\mathrm{d}p}{\mathrm{d}x},
\]
其中扩散系数与迁移率由爱因斯坦关系联系：
\[
D_n=\mu_nU_T,\qquad D_p=\mu_pU_T .
\]
代入硅的数值，$D_n=1350\times0.0259\approx34.9\,\mathrm{cm^2/s}$。
若界面附近 $1\,\mathrm{\mu m}$ 内浓度变化 $10^{15}\,\mathrm{cm^{-3}}$，则电子扩散电流密度
\[
J_n=qD_n\frac{\Delta n}{\Delta x}
=1.6\times10^{-19}\times34.9\times\frac{10^{15}}{10^{-4}}
\approx5.6\times10^{-2}\,\mathrm{A/cm^2}.
\]
这个量级说明：\textbf{只要浓度差还在，扩散就是一股很大的电流}。

p34 接着说，越过界面的载流子留下了不能移动的固定离子：
「As they cross the junction, the fixed ions they leave behind establish a 「built-in」 electric
field at the junction.」。N 侧留下正离子、P 侧留下负离子，于是结区出现一层
「N 侧正、P 侧负」的空间电荷，电场方向由 N 指向 P。

\textbf{空间电荷的分配}：结区两侧的电荷量必须相等（电中性），即
\[
qN_D x_n=qN_A x_p,\qquad x_n+x_p=W .
\]
解出
\[
x_n=\frac{N_A}{N_A+N_D}W,\qquad x_p=\frac{N_D}{N_A+N_D}W .
\]
耗尽层\textbf{主要向轻掺杂一侧扩展}。举一个不对称的例子：$N_A=10^{17}$、$N_D=10^{15}\,\mathrm{cm^{-3}}$，
$V_{bi}\approx0.693\,\mathrm{V}$、$W\approx0.95\,\mu\mathrm{m}$，则 $x_n\approx0.94\,\mu\mathrm{m}$、$x_p\approx0.009\,\mu\mathrm{m}$——
99\% 的耗尽层落在轻掺杂的 N 侧。单位面积空间电荷
\[
Q=qN_D x_n\approx1.6\times10^{-19}\times10^{15}\times0.94\times10^{-4}
\approx1.5\times10^{-8}\,\mathrm{C/cm^2}.
\]
这个 $Q$ 是后面结电容与击穿计算的出发点。

\subhdr{3.3 漂移、平衡与内建电势：$0.7\,\mathrm{V}$ 的真正来源（p35--p36）}

p35 的定义是 「Drift motion: the movement caused by electric field」。内建电场把电子往 N 侧推、
把空穴往 P 侧推——\textbf{方向与扩散相反}。随着扩散进行，空间电荷越来越多、电场越来越强，
漂移电流也越来越大，直到两者相等：
\[
J_n=q n\mu_n E+qD_n\frac{\mathrm{d}n}{\mathrm{d}x}=0 .
\]
这就是 p36 说的两种效应并存（「At the junction there are two effects occurring」）与动态平衡。
注意「平衡」不等于「静止」：两个方向各有很大的电流，只是净电流为零。

把 $E=-\mathrm{d}\varphi/\mathrm{d}x$、$D_n=\mu_nU_T$ 代入上式并整理：
\[
-n\frac{\mathrm{d}\varphi}{\mathrm{d}x}+U_T\frac{\mathrm{d}n}{\mathrm{d}x}=0
\quad\Longrightarrow\quad
\mathrm{d}\varphi=U_T\frac{\mathrm{d}n}{n}.
\]
从 P 区积分到 N 区，用两侧的平衡多子浓度 $n_p=n_i^2/N_A$、$n_n\approx N_D$：
\[
V_{bi}=\varphi_n-\varphi_p
=U_T\ln\frac{n_n}{n_p}
=U_T\ln\frac{N_A N_D}{n_i^2}.
\]
\begin{center}
\begin{minipage}{0.94\linewidth}
\subhdr{内建电势（原形式与数值）}
设突变 PN 结两侧的受主、施主浓度分别为 $N_A,N_D$，温度 $T$，本征浓度 $n_i$，则平衡时的内建电势为
\[
V_{bi}=U_T\ln\frac{N_A N_D}{n_i^2},\qquad U_T=\frac{kT}{q}.
\]
取 $N_A=N_D=10^{16}\,\mathrm{cm^{-3}}$、300 K：
\[
V_{bi}(\mathrm{Si})=0.0259\ln\frac{10^{32}}{2.25\times10^{20}}\approx0.693\,\mathrm{V},
\]
\[
V_{bi}(\mathrm{Ge})=0.0259\ln\frac{10^{32}}{5.76\times10^{26}}\approx0.312\,\mathrm{V}.
\]
\aside{课件 p40 的 「In silicon, this occurs at about 0.7 volts forward bias」 与 p52 的
$U_D=0.7\,\mathrm{V}$(Si)、$0.3\,\mathrm{V}$(Ge)，数值正对应上面两个 $V_{bi}$。
但课件从未写出这两个数的来源——本讲义这一节就是它的出处。}
\end{minipage}
\end{center}

\textbf{换一条路再算一遍（检验）}：平衡时，同一批电子在 N 区与 P 区的概率比由玻尔兹曼因子给出，
\[
\frac{n_n}{n_p}=e^{qV_{bi}/(kT)}=e^{V_{bi}/U_T}
\quad\Longrightarrow\quad
V_{bi}=U_T\ln\frac{n_n}{n_p},
\]
与积分结果完全一致。两条路一条从「净电流为零」出发、一条从「统计分布」出发，
结论相同，说明 $V_{bi}$ 是热平衡的必然，不是某个模型的假设。

\textbf{这里必须停一下，纠正一个常见误解}：$0.7\,\mathrm{V}$ 不是「外加电压把势垒恰好抵消」的开关阈值，
而是「势垒被压到接近零、多数载流子扩散电流达到 mA 量级」的一个工程约定。
严格说电流从 $0.5\,\mathrm{V}$ 起就按指数上升（§5.2 给公式），
$0.7\,\mathrm{V}$ 只是这条指数曲线在 mA 量级上的「代表点」。
课件 p40 用 「about」 是准确的，把 $0.7\,\mathrm{V}$ 当普适常数就错了（§6.2 再展开）。

\subhdr{3.4 耗尽层有多宽、电场有多强}

平衡时空间电荷全部集中在界面两侧，形成一个几乎没有载流子的耗尽层。由泊松方程解得（突变结）
\[
W=\sqrt{\frac{2\varepsilon_s\left(V_{bi}-U\right)}{q}\left(\frac{1}{N_A}+\frac{1}{N_D}\right)},
\]
其中 $\varepsilon_s=11.7\varepsilon_0$（硅）、$U$ 是外加电压（正偏取正、反偏取负）。
代入 $N_A=N_D=10^{16}\,\mathrm{cm^{-3}}$、$U=0$：
\[
W=\sqrt{\frac{2\times1.035\times10^{-12}\times0.693}{1.6\times10^{-19}}\times2\times10^{-16}}
\approx4.2\times10^{-5}\,\mathrm{cm}=0.42\,\mu\mathrm{m}.
\]
最大电场在界面处，$E_{\max}=2(V_{bi}-U)/W\approx3.3\times10^{4}\,\mathrm{V/cm}$。
结电容（势垒电容）$C_T=\varepsilon_s A/W$，随反偏电压增大而减小——这是变容二极管的原理，
也是 §7 讲参数时「结电容」一项的来处。[教材补充：$W$、$E_{\max}$、$C_T$ 三条课件未写，
补它是因为后面解释「正偏变薄、反偏变厚」需要宽度这个量。]

\subhdr{3.5 平衡的定量图像}

把三个量放在一起看：$V_{bi}\approx0.69\,\mathrm{V}$、$W\approx0.42\,\mu\mathrm{m}$、
$E_{\max}\approx3.3\times10^{4}\,\mathrm{V/cm}$。数量级上
$V_{bi}\approx E_{\max}W/2$ 完全自洽。这三个数在 §4 会各用一次：
$V_{bi}$ 决定导通阈值，$W$ 决定正/反偏时「变薄还是变厚」，$E_{\max}$ 决定击穿时场强是否够高。

%=============================================================
\subhdr{3.6 把势垒画在能带上：三个名词的关系}

前面用「势垒高度 $qV_{bi}$」描述平衡，用「电场」描述方向。这两个说法在能带图上是一回事：
N 区电子多、费米能级靠近导带；P 区空穴多、费米能级靠近价带。两块接触后费米能级必须拉平，
于是能带在结区发生弯曲，弯曲量正好是 $qV_{bi}$：
\[
E_{C}(p)-E_{C}(n)=qV_{bi}.
\]
三句话记牢：
（1）\textbf{能带弯曲量 = 势垒高度}，数值 $0.693\,\mathrm{V}$（硅，$10^{16}$ 掺杂）；
（2）\textbf{正偏 $U$ 使弯曲变平}，剩余势垒 $q(V_{bi}-U)$，耗尽层变薄（p40 的 「thinner」）；
（3）\textbf{反偏使弯曲更陡}，势垒 $q(V_{bi}+|U|)$，耗尽层变厚（p42 的 「thicker」）。
能带图的价值在于：它一眼看出「为什么正偏能导通、反偏不能」——不是载流子变多了，
而是同一个势垒被外电压抬高或压低。
\aside{p39 说两个场「相加」，能带图上就是两条 $E_C$ 曲线的斜率相加或相减；
p42 的「同向」就是两段弯曲同向叠加。}

\subhdr{3.7 结电容：反偏为什么能当变容管}

耗尽层两侧的电荷量随电压变，所以 PN 结本身就是一个电容。由 §3.2 的电荷分配，
单位面积电荷
\[
Q(U)=\sqrt{2q\varepsilon_s\left(V_{bi}-U\right)\frac{N_AN_D}{N_A+N_D}},
\]
对电压求导得微分电容
\[
C_T=\left|\frac{\mathrm{d}Q}{\mathrm{d}U}\right|=\frac{\varepsilon_s}{W},
\]
乘上结面积 $A$ 得 $C_T=\varepsilon_sA/W$。因为 $W\propto\sqrt{V_{bi}-U}$，所以
$C_T\propto1/\sqrt{V_{bi}-U}$：反偏越大，$W$ 越宽，电容越小。
以 $N_A=10^{17}$、$N_D=10^{15}\,\mathrm{cm^{-3}}$、$A=1\,\mathrm{mm^2}$ 为例：
$U=0$ 时 $W=0.95\,\mu\mathrm{m}$、$C_T\approx109\,\mathrm{pF}$；
$U=-5\,\mathrm{V}$ 时 $W=2.73\,\mu\mathrm{m}$、$C_T\approx38\,\mathrm{pF}$，减小到三分之一。
这就是变容二极管调谐的原理，也是高频电路里「结电容」必须计入的原因。
\aside{正偏时 $V_{bi}-U$ 变小，$W$ 变薄，$C_T$ 变大，但正偏同时有大量扩散电容，
不能再用这个公式；本节只谈反偏。}

\hdr{§4}{单向导电：加电压时结里发生了什么}

\subhdr{4.1 正偏：外场与内场反向，势垒被压低（p37--p40）}

p37 的接法是「P 侧接电源正极、N 侧接负极」（图 2：low 在 N、High 在 P），这叫正向偏置。
p39 的原句是 「At the P-N junction, the 「built-in」 electric field and the applied electric field
are in opposite directions. When these two fields add, the resultant field at the junction is
smaller in magnitude than the magnitude of the original 「built-in」 electric field.」。

「两个场相加」这句话需要写成算式。设外加电压为 $U$（P 侧相对 N 侧为正，$U>0$），
外电场方向由 P 指向 N，与内建电场（N$\to$P）相反，于是净势垒高度变成
\[
q\left(V_{bi}-U\right).
\]
$U$ 越大，势垒越低，能翻过势垒的多子越多，扩散电流越大。p40 的结论是
「This results in a thinner, less resistive depletion region. If the applied voltage is large
enough, the depletion region's resistance becomes negligible. In silicon, this occurs at about
0.7 volts forward bias.」。把它与 §3.4 的 $W$ 公式对照：$V_{bi}-U$ 变小，$W$ 确实变薄。

\subhdr{4.2 反偏：势垒抬高，只剩少子漂移（p41--p42）}

p41 的接法反过来：正电压加在 N 侧、负电压加在 P 侧（图 3：High 在 N、Low 在 P）。
此时外电场与内建电场\textbf{同向}，p42 的原句是 「The 「built-in」 electric field and the applied
electric field are in the same direction. When these two fields add, the resultant larger electric
field is in the same direction as the 「built in」 electric field and this creates a thicker,
more resistive depletion region.」。

势垒变成 $q(V_{bi}+|U|)$，多子几乎翻不过去，扩散电流趋近于零；剩下的只有
\textbf{少子}被电场拉过结的漂移电流。少子浓度由温度决定、与外加电压几乎无关，
所以这段电流近似是一个常数——这就是 p49 标注的 「Reverse saturation current」。

反偏电压再加大，结内场强升高到某个值时会发生击穿（p49 的 「Reverse breakdown」，
横轴标 $U_{BR}$）。两种机制：
（1）齐纳击穿：掺杂很重、耗尽层很薄，场强足以直接把价带电子隧道到导带，一般发生在 $U_{BR}<5\,\mathrm{V}$；
（2）雪崩击穿：载流子在强场中加速，碰撞电离产生新的电子—空穴对，连锁倍增，一般发生在 $U_{BR}>7\,\mathrm{V}$。
击穿本身不毁管子，只要外电路把电流限制住；但普通整流管不允许长期工作在击穿区。
[教材补充：两种击穿的区分课件未写，p49 只标了 $U_{BR}$。]

\subhdr{4.3 ON/OFF 的判据（p43）}

p43 用两张双色方块图作总结：正偏时 N 侧 low、P 侧 High，标 \textbf{ON}；
反偏时 N 侧 High、P 侧 Low，标 \textbf{OFF}。把它写成可执行的判据：

\begin{center}
\begin{minipage}{0.94\linewidth}
\callout{\textbf{单向导电判据}\quad 记 $U_{PN}=\varphi_P-\varphi_N$。\\
（1）$U_{PN}>U_D$（硅约 $0.7\,\mathrm{V}$）：导通，相当于开关闭合；\\
（2）$U_{PN}<U_D$：截止，相当于开关断开；\\
（3）反向（$U_{PN}<0$）：截止，只有 $\mu\mathrm{A}$ 级以下的反向电流；\\
（4）反向电压达到 $U_{BR}$：击穿，普通管不允许。}
\end{minipage}
\end{center}

这里最容易错的是把「正偏」理解成「P 侧接正极」就完了。真正的条件是
\textbf{P 侧电位比 N 侧高出一个导通压降}。若 P 侧接正极但只高出 $0.2\,\mathrm{V}$，
硅管仍然截止——p57 的钳位题里，D\_A 的阳极电位只有 $0.7\,\mathrm{V}$、阴极 $3\,\mathrm{V}$，
虽然「方向没接反」，它依然截止。§8 会反复用到这一点。

\subhdr{4.4 为什么反向电流几乎不随电压变}

反向时能翻过结的只有少子，而少子的数量由\textbf{热激发}决定，与电压无关；
电压只改变它们被拉过结的速度，不改变供给量。因此反向电流像是一个与电压无关的电流源 $I_S$。
温度每升高约 $10\,{}^\circ\mathrm{C}$，$n_i^2$ 增大约 4 倍，$I_S$ 也随之增大——
这解释了 p53--p54 那张「高温曲线反向段下沉更深」的图（§7.1 定量）。

%=============================================================
\subhdr{4.5 击穿的定量图像与稳压管}

p49 只标了 $U_{BR}$，没给击穿的数学描述。工程上用一个倍增因子 $M$ 描述雪崩：
\[
M=\frac{1}{1-\left(U/U_{BR}\right)^{m}},
\]
其中 $m=3\sim6$。$U\to U_{BR}$ 时 $M\to\infty$，电流失控；所以「击穿」不是电压到此为止，
而是\textbf{电流不再受电压控制}。只要外电路有限流电阻，管子不会坏；这就是稳压管的工作方式。

\textbf{齐纳击穿与雪崩击穿的分界与温度系数}：
（1）$U_{BR}<5\,\mathrm{V}$：耗尽层很薄，隧道效应为主（齐纳），温度升高时 $U_{BR}$ 反而下降；
（2）$U_{BR}>7\,\mathrm{V}$：碰撞电离为主（雪崩），温度升高时 $U_{BR}$ 上升；
（3）$5$--$7\,\mathrm{V}$：两种机制混合，温度系数接近零——所以稳压管常选 $6\,\mathrm{V}$ 左右。
[教材补充：倍增因子与温度系数课件未写；p56 的参数表只列了 $U_{BR}$。]

稳压管与普通整流管的区别只有一条：\textbf{它被设计成长期工作在击穿区}，
所以参数里多了稳定电压 $U_Z$、动态电阻 $r_Z$、最小与最大稳定电流 $I_{Z\min},I_{Z\max}$。
普通管一旦进入击穿区，即使电流受限，长期工作也会因结温过高而损坏。

\subhdr{4.6 反向饱和电流的温度系数}

§5.2 会给出 $I_S\propto n_i^2\propto T^3e^{-E_g/(kT)}$。取对数再对温度求导，
\[
\frac{1}{I_S}\frac{\mathrm{d}I_S}{\mathrm{d}T}
=\frac{3}{T}+\frac{E_g}{kT^2}.
\]
代入硅的数值（$T=300\,\mathrm{K}$、$E_g=1.12\,\mathrm{eV}$）：
\[
\frac{3}{300}+\frac{1.12}{8.617\times10^{-5}\times300^2}
\approx0.01+0.144\approx0.15\ \mathrm{K^{-1}}.
\]
即温度每升 $1\,{}^\circ\mathrm{C}$，理想反向电流增大约 $15\%$，每升 $10\,{}^\circ\mathrm{C}$ 约增至 4 倍；
从 $20$ 到 $80\,{}^\circ\mathrm{C}$ 约增大 $3.3\times10^3$ 倍。
实测值因表面漏电、复合电流等机制通常没有这么快（常见口径是每 $10\,{}^\circ\mathrm{C}$ 翻倍），
但「指数增长」这个结论不变：\textbf{反向漏电是温度的强函数}，高温选管必须按最高结温降额。

\subhdr{4.7 两种电容：结电容与扩散电容}

§3.7 的 $C_T$ 是\textbf{反偏}时的主角；正偏时还有另一个更大的电容——扩散电容 $C_D$。
它来自中性区存储的少子电荷：正偏时大量少子被注入并堆积，电压一变，堆积量跟着变，
等效出一个电容
\[
C_D=\frac{\tau I_D}{nU_T},
\]
其中 $\tau$ 是少子寿命。取 $\tau=100\,\mathrm{ns}$、$I_D=1\,\mathrm{mA}$、$n=1.3$，
\[
C_D\approx\frac{10^{-7}\times10^{-3}}{1.3\times0.0259}\approx3.0\,\mathrm{nF},
\]
比反偏时的 $38\,\mathrm{pF}$（§3.7）大近两个数量级。

\textbf{为什么重要}：二极管从正偏切到反偏时，先要把存储的电荷抽走，这段「反向恢复时间」
限制了开关速度。p48 把开关二极管单列一类，就是为了这个指标。
[教材补充：扩散电容与反向恢复不在课件中，补它是为了回答「正偏时结电容为什么不能用 $C_T$」。]

\hdr{§5}{二极管的伏安特性：五个名词各指哪一段}

\subhdr{5.1 结构、符号与封装（p44--p48）}

p46 给定结构：「A single PN junction with contacts for connecting to external circuits」——
一个 PN 结加两条引线。P 侧引出的是阳极（Anode），N 侧是阴极（Cathode）（p47）。
电路符号是一个实心三角形加一条竖线：\textbf{三角形由阳极指向阴极}，竖线在阴极一侧。
做题时先看三角形朝哪，正向电流就只能朝那个方向流。

p46 还给了三条实用信息：塑封管的白环一端是负极、金属封装的螺栓端是负极、玻璃封管的黑环是负极。
这三条在接实物时才用得上，考试不考，但装电路时接反会烧管子。

p48 的分类是三条并列标准：（1）按结构：点接触型、面结型；（2）按材料：硅管、锗管；
（3）按用途：整流管、开关管。这三条分别对应「结电容大小」「$0.7$ 还是 $0.3$」「能过多大电流」，
是后面选参数的索引。

\subhdr{5.2 曲线的四个区与五个名词（p49--p50）}

p49 画了硅管（实线）与锗管（虚线）的伏安特性；p50 把同一条曲线配上五个问句，
\textbf{一个答案都没给}。这五个问句就是本节要回答的全部内容：

\par\vspace{2pt}
\noindent\begin{tabular}{@{}>{\raggedright\arraybackslash}p{6.9cm}>{\raggedright\arraybackslash}p{8.9cm}@{}}
\tabtop
\thead{课件的问句（p50）} & \thead{答案}
\\
\tabmid
(1) Which part is forward biased? & 第一象限 $U>0$ 的整段，从原点到曲线陡升段
\\
(2) Which part is reverse biased? & 第三象限 $U<0$ 的整段，曲线贴近横轴
\\
(3) What is Barrier potential? & 正向曲线开始明显上翘的电压，硅约 $0.5$--$0.7\,\mathrm{V}$、锗约 $0.2$--$0.3\,\mathrm{V}$
\\
(4) What is reverse saturation current? & 反向截止段那个几乎与电压无关的电流 $I_R$，硅 $\mathrm{nA}$ 级、锗 $\mu\mathrm{A}$ 级
\\
(5) What is reverse breakdown voltage? & 反向曲线向下急折处的电压 $U_{BR}$
\tabbot
\end{tabular}

曲线为什么长这样，课件没有给方程。把 §3.3 的势垒图像写成电流—电压关系，
就是下面这条式子（肖克利方程）：

\begin{center}
\begin{minipage}{0.96\linewidth}
\subhdr{肖克利方程（本讲义补的推导）}
设中性区少子满足扩散方程，边界处少子浓度由玻尔兹曼关系给出
\[
\frac{\mathrm{d}^2p_n}{\mathrm{d}x^2}=\frac{p_n-p_{n0}}{L_p^2},
\qquad L_p=\sqrt{D_p\tau_p},
\]
\[
p_n(x)=p_{n0}+\left[p_n(0)-p_{n0}\right]e^{-x/L_p},
\qquad p_n(0)=p_{n0}\,e^{U/U_T}.
\]
在 $x=0$ 处取扩散电流，
\[
I_p(0)=qA D_p\left.\frac{\mathrm{d}p_n}{\mathrm{d}x}\right|_{x=0}
=\frac{qA D_p p_{n0}}{L_p}\left(e^{U/U_T}-1\right).
\]
电子一侧同理，两项相加得
\[
I=I_S\left(e^{U/(nU_T)}-1\right),
\qquad
I_S=qA\left(\frac{D_p p_{n0}}{L_p}+\frac{D_n n_{p0}}{L_n}\right).
\]
其中 $n$ 是理想因子（硅管 $1\sim2$，本讲义取 $1.3$ 做数值例）。
\aside{课件 p49--p50 只画曲线、不给方程，本式是 §4.1「势垒越低电流越大」的定量版；
$p_{n0}=n_i^2/N_D$、$n_{p0}=n_i^2/N_A$，所以 $I_S$ 与 $n_i^2$ 成正比——这条关系在 §7 的温度题里要用。}
\end{minipage}
\end{center}

由这条式子可以直接读出四个区：
（1）$U\gg nU_T$ 时 $e^{U/(nU_T)}\gg1$，电流按指数上升——这就是「导通」；
（2）$U$ 在 $0$ 到 $0.5\,\mathrm{V}$ 之间时指数还很小，宏观上看像「死区」；
（3）$U<0$ 且 $|U|\gg nU_T$ 时 $e^{U/(nU_T)}\to0$，$I\to-I_S$，与电压无关——反向饱和；
（4）$|U|$ 大到击穿时，上式失效，把四个区写成三段可算的形式。正向导通时
\[
I\approx I_S\,e^{U/(nU_T)}\qquad(U\gg nU_T);
\]
反向截止时
\[
I\approx-I_S\qquad(U<0,\ |U|\gg nU_T);
\]
死区里电流小到可以当零：
\[
|I|\lesssim I_S\left(e^{0.5/(nU_T)}-1\right)\qquad(0<U<0.5\,\mathrm{V}).
\]
以 $I_S=10^{-12}\,\mathrm{A}$、$n=1.3$ 代入，$U=0.5\,\mathrm{V}$ 时只有约 $3\,\mu\mathrm{A}$，
比 mA 级的导通电流小三个数量级——这就是「死区」可以当零的定量依据。

进入击穿机制。

\subhdr{5.3 反向饱和电流从哪来、有多小}

$I_S$ 的表达式里两个括号项都正比于 $n_i^2$，而 $n_i^2\propto e^{-E_g/(kT)}$。
硅的禁带 $1.12\,\mathrm{eV}$ 比锗的 $0.67\,\mathrm{eV}$ 宽，所以同样温度下硅的 $I_S$ 比锗小几个数量级：
硅典型 $10^{-12}$--$10^{-9}\,\mathrm{A}$，锗典型 $10^{-7}$--$10^{-6}\,\mathrm{A}$。
这也解释了 p49 图上锗管反向段比硅管低得多的画法。

\subhdr{5.4 硅管与锗管的曲线对比（p49）}

p49 的两条曲线差在三处：
（1）拐点位置：硅在 $0.5$--$0.7\,\mathrm{V}$，锗在 $0.2$--$0.3\,\mathrm{V}$（因为 $V_{bi}$ 不同，§3.3）；
（2）反向电流量级：锗大得多（因为 $n_i$ 大得多）；
（3）同样电流下的管压降：硅高、锗低，所以锗管功耗小但漏电大、温度稳定性差。
这三条正是「硅管为什么成了主流」的答案。

\subhdr{5.5 用 p55 的两张图反推材料}

p55 给了两张量程不同的伏安曲线，问 「which one is for Si-diode and which one is for Ge-diode?」。
判别只看拐点：拐点在 $0.6$--$0.7\,\mathrm{V}$ 附近（纵轴 $20/40/60\,\mathrm{mA}$ 那张）是硅管；
拐点在 $0.2$--$0.3\,\mathrm{V}$ 附近（纵轴 $5/10/15\,\mathrm{mA}$ 那张）是锗管。
[待确认：扫描页的左右排列与文字层抽取顺序相反，本讲义按「拐点电压」判别，不按左右位置记忆；
若考试给的原图与这里相反，以拐点为准。]

\textbf{这里要指出课件的一处口径不一致}：p52 写锗管 $U_D=0.3\,\mathrm{V}$，
而 p55 的锗管曲线拐点读出来约 $0.2\,\mathrm{V}$。两者不是矛盾，而是「读图位置」不同：
$0.3\,\mathrm{V}$ 是 mA 级电流下的工程约定值，曲线拐点是电流刚开始抬头的位置。
本讲义统一按 $0.3\,\mathrm{V}$ 计算，读图时说明「约 $0.2$--$0.3\,\mathrm{V}$」。

%=============================================================
\subhdr{5.6 读图三练：从曲线上把三个数取下来}

p50 的五个问句本质上是「在同一张图上读三个数」。读图的规矩是\textbf{先定坐标、再定区域、最后读交点}：
（1）读死区电压：沿横轴找曲线开始离开横轴的位置。p55 两张图分别在 $0.6$--$0.7\,\mathrm{V}$ 与
$0.2$--$0.3\,\mathrm{V}$ 一带起翘，前者约为 $0.6\,\mathrm{V}$、后者约为 $0.2\,\mathrm{V}$
（左右顺序的说明见 §5.5）；
（2）读反向饱和电流：沿纵轴负方向找那段水平线的纵坐标。p49 的硅管在 $\mathrm{nA}$ 级、锗管在 $\mu\mathrm{A}$ 级，
图上的 $-0.02\,\mathrm{mA}$ 与 $-0.04\,\mathrm{mA}$ 刻度是为锗管留的；
（3）读击穿电压：找曲线向下急折处的横坐标。p49/p55 的 $U_{BR}$ 在 $-50\,\mathrm{V}$ 附近。
三者读错任何一个，后面选参数都会错：死区电压决定导通判断，$I_R$ 决定漏电功耗，$U_{BR}$ 决定耐压。

\subhdr{5.7 一条曲线同时读出温度与材料}

p53 与 p55 的两组题，本质是同一个判据的两次应用：
（1）\textbf{同材料、比温度}：拐点更靠左、反向电流更大的是高温曲线（p54 的红色 $80\,{}^\circ\mathrm{C}$）；
（2）\textbf{同温度、比材料}：拐点更靠左、反向电流更大的是锗管（p55）。
两者的现象方向相同，但原因不同：温度改变的是 $n_i(T)$，材料改变的是 $E_g$。
判据可以合成一句话：\textbf{看拐点定电压、看反向段定量级、再看题目问的是温度还是材料}。

\subhdr{5.8 四个参数在曲线上的位置}

p56 把 $I_F$、$U_F$、$U_R$、$I_{Rm}$ 直接标在一条曲线上，但没说明各自「量的是哪一段」。
补一张对照：

\par\vspace{2pt}
\noindent\begin{tabular}{@{}>{\raggedright\arraybackslash}p{2.0cm}>{\raggedright\arraybackslash}p{3.0cm}>{\raggedright\arraybackslash}p{4.6cm}>{\raggedright\arraybackslash}p{5.0cm}@{}}
\tabtop
\thead{符号} & \thead{名称} & \thead{在曲线的哪一段} & \thead{怎么读}
\\
\tabmid
$I_F$ & 额定正向平均电流 & 正向导通段的工作点 & 不是曲线上一点，由结面积与散热决定
\\
$U_F$ & 正向压降 & 正向曲线在 $I_F$ 处的横坐标 & 同一只管子电流越大 $U_F$ 越高
\\
$U_R$ & 最高反向工作电压 & 反向平坦段 & 取 $U_{BR}$ 的 $1/2$--$2/3$，不是 $U_{BR}$ 本身
\\
$I_{Rm}$ & 最大反向电流 & 反向平坦段的纵坐标 & 室温值，高温要放大
\\
$U_{BR}$ & 反向击穿电压 & 反向急折点 & 只用于选耐压，不允许长期工作
\tabbot
\end{tabular}

\textbf{读表的关键}：$I_F$ 与 $I_{Rm}$ 是\textbf{极限值}，$U_F$ 与 $U_{BR}$ 是\textbf{特性值}。
极限值用来「校核」，特性值用来「计算」——两者混用是选型题最常见的错。

\hdr{§6}{三种等效模型：什么时候可以把它当开关}

课件只给了两种模型（p51 理想、p52 恒压降），而实际做题需要三种。
第三种（折线模型）标 [教材补充]，补它的理由写在 §6.3 开头。

\subhdr{6.1 理想模型（p51）}

p51 的写法是：「Forward bias\quad $U_D=0$，$R_D=0$\quad short circuit」，
「Reverse bias\quad $I_D=0$，$R_D\to\infty$\quad open circuit」。
也就是把二极管当成一个\textbf{理想开关}：正向短路、反向开路，两端压降恒为零。

它什么时候能用：电源电压远大于 $0.7\,\mathrm{V}$（比如 $12\,\mathrm{V}$ 系统里判断逻辑电平），
或者题目只问「导通/截止」而不问具体电压。误差量级是 $0.7/U_S$：
$U_S=12\,\mathrm{V}$ 时约 $6\%$，$U_S=1\,\mathrm{V}$ 时就到了 $70\%$，不能再用。

\subhdr{6.2 恒压降模型（p52）}

p52 的写法是「$U_D=0.7\,\mathrm{V}$(Si)，$U_D=0.3\,\mathrm{V}$(Ge)」，反向仍开路。
正向时把它等效成一个恒压源 $U_D$ 串联在回路里，也就是「先扣掉 $0.7\,\mathrm{V}$，剩下的电压才加在电阻上」。

\textbf{$0.7\,\mathrm{V}$ 从哪来}：§3.3 算出硅的 $V_{bi}\approx0.693\,\mathrm{V}$，
数值上正是 $0.7\,\mathrm{V}$；锗的 $V_{bi}\approx0.312\,\mathrm{V}$，正是 $0.3\,\mathrm{V}$。
所以「势垒电压」这个叫法不是比喻，它确实是内建电势。

但\textbf{它不是一个常数}。用 §5.2 的肖克利方程（取 $I_S=10^{-12}\,\mathrm{A}$、$n=1.3$）算同样一只硅管：
\[
U=nU_T\ln\left(\frac{I}{I_S}+1\right),
\]
\[
I=0.1\,\mathrm{mA}\Rightarrow U\approx0.619\,\mathrm{V},\quad
I=1\,\mathrm{mA}\Rightarrow U\approx0.696\,\mathrm{V},\quad
I=10\,\mathrm{mA}\Rightarrow U\approx0.774\,\mathrm{V}.
\]
电流变 100 倍，压降才变 $0.15\,\mathrm{V}$——这就是「恒压降」三个字的底气，也是它的误差边界：
mA 级、室温下取 $0.7\,\mathrm{V}$ 误差在 $0.1\,\mathrm{V}$ 以内；电流低到 $\mu\mathrm{A}$ 级或温度变化几十度，
就必须回到指数式（§7.1）。

\subhdr{6.3 折线模型（第三种，[教材补充]）}

课件没给第三种模型。补它是因为两种模型都假设「导通后压降不变」，
而真实管子在电流大时还有一个随电流增长的\emph{体电阻}压降，在小信号分析里还要用\emph{动态电阻}。
折线模型把正向段写成一条斜线：
\[
U_D=U_{th}+I_D r_d,\qquad r_d=\frac{nU_T}{I_D},
\]
其中 $U_{th}$ 是折线的横轴截距（硅取 $0.6$--$0.7\,\mathrm{V}$），$r_d$ 是工作点处的动态电阻。
$r_d$ 的量级很说明问题：
\[
I_D=0.1\,\mathrm{mA}\Rightarrow r_d\approx336\,\Omega,\quad
1\,\mathrm{mA}\Rightarrow33.6\,\Omega,\quad
10\,\mathrm{mA}\Rightarrow3.4\,\Omega,\quad
100\,\mathrm{mA}\Rightarrow0.34\,\Omega.
\]
电流越小，动态电阻越大——所以「小信号」分析里不能把二极管当恒压源，要当电阻。

\subhdr{6.4 同一道题、三种模型：误差到底多大}

取一个最简单的电路：$U_S=10\,\mathrm{V}$ 的直流源串联 $R=1\,\mathrm{k\Omega}$ 再串一只硅二极管，
求 $I_D$ 与 $U_D$。三种模型与「精确解」（肖克利方程，$I_S=10^{-12}\,\mathrm{A}$、$n=1.3$）对比如下：

\par\vspace{2pt}
\noindent\begin{tabular}{@{}>{\raggedright\arraybackslash}p{2.3cm}>{\raggedright\arraybackslash}p{3.4cm}>{\raggedright\arraybackslash}p{2.5cm}>{\raggedright\arraybackslash}p{2.6cm}>{\raggedright\arraybackslash}p{3.0cm}@{}}
\tabtop
\thead{模型} & \thead{方程} & \thead{$I_D$} & \thead{$U_D$} & \thead{相对误差}
\\
\tabmid
精确解 & $I=I_S(e^{U/nU_T}-1)$ & $9.229\,\mathrm{mA}$ & $0.771\,\mathrm{V}$ & 基准
\\
理想 & $U_D=0$ & $10.00\,\mathrm{mA}$ & $0$ & $I$ 偏大 $8.4\%$
\\
恒压降 & $U_D=0.7\,\mathrm{V}$ & $9.300\,\mathrm{mA}$ & $0.700\,\mathrm{V}$ & $I$ 偏大 $0.8\%$、$U$ 偏小 $9.2\%$
\\
折线 & $U_D=0.7+3.6I_D$ & $9.267\,\mathrm{mA}$ & $0.734\,\mathrm{V}$ & $I$ 偏大 $0.4\%$、$U$ 偏小 $4.9\%$
\tabbot
\end{tabular}

\textbf{读表三条}：
（1）只要电源电压是几伏、电阻是千欧级，恒压降模型的电流误差不到 $1\%$——这就是它成为手算主力的原因；
（2）理想模型的电压误差是 $100\%$（它把 $0.77\,\mathrm{V}$ 当成 $0$），只在「判断通断」时可以接受；
（3）折线模型把电压误差从 $9.2\%$ 压到 $4.9\%$，代价是多一个参数 $r_d$。
选哪个模型，取决于题目问的是「通断」「电流」还是「电压」。

\par\vspace{2pt}
\noindent\begin{tabular}{@{}>{\raggedright\arraybackslash}p{2.2cm}>{\raggedright\arraybackslash}p{3.4cm}>{\raggedright\arraybackslash}p{2.6cm}>{\raggedright\arraybackslash}p{4.4cm}>{\raggedright\arraybackslash}p{2.4cm}@{}}
\tabtop
\thead{模型} & \thead{正向等效} & \thead{反向} & \thead{适用场合} & \thead{参数个数}
\\
\tabmid
理想 & 短路（$U_D=0$） & 开路 & 电源 $\gg0.7\,\mathrm{V}$、只判通断、逻辑电平 & 0
\\
恒压降 & $U_D$ 恒压源 & 开路 & 几伏以上整流/钳位的手算，mA 级 & 1
\\
折线 & $U_{th}+I r_d$ & 开路（可加漏电） & 大电流、精密计算、小信号动态电阻 & 2
\tabbot
\end{tabular}

%=============================================================
\textbf{再补一个反向的算例}：同一只管子反接在 $10\,\mathrm{V}$ 电源上（$U_D=-10\,\mathrm{V}$），
串 $1\,\mathrm{k\Omega}$。用肖克利方程
\[
I=I_S\left(e^{-10/(nU_T)}-1\right)\approx-I_S=-10^{-12}\,\mathrm{A},
\]
电阻上的压降只有 $10^{-9}\,\mathrm{V}$——\textbf{反向电流几乎全部由管子自己决定，与串联电阻无关}。
这就是「反向截止」在数值上的样子：$10\,\mathrm{V}$ 几乎全落在管子上，回路电流约 $1\,\mathrm{pA}$。
理想模型说「开路」、恒压降模型也说「开路」，两者在这一步没有差别；差别只在正向。

\subhdr{6.5 小信号：把二极管在 Q 点换成一个电阻}

前四个小节处理的都是「直流大信号」。若二极管上叠加一个很小的交流信号 $\Delta u$，
它不再是一个恒压源，而是在工作点 $Q$ 附近的一条斜线，斜率就是动态电阻：
\[
r_d=\left.\frac{\mathrm{d}U_D}{\mathrm{d}I_D}\right|_{Q}
=\frac{nU_T}{I_D}.
\]
取 $n=1.3$、$I_D=1\,\mathrm{mA}$，$r_d\approx33.6\,\Omega$。
若叠加 $\Delta u=1\,\mathrm{mV}$，线性化给出
\[
\Delta i=\frac{\Delta u}{r_d}=\frac{1\,\mathrm{mV}}{33.6\,\Omega}\approx30\,\mu\mathrm{A},
\]
而按指数式精确计算为 $30.2\,\mu\mathrm{A}$，误差约 $1\%$——线性化在 $\Delta u\ll nU_T$ 时成立。
若把 $\Delta u$ 加到 $10\,\mathrm{mV}$（已与 $nU_T=33.6\,\mathrm{mV}$ 同量级），
线性化给出 $298\,\mu\mathrm{A}$、精确值 $347\,\mu\mathrm{A}$，误差约 $14\%$——这就是「小信号」三个字的边界。
\aside{本册只讲直流，但下一册的微变等效电路（C4-3）正是从这里开始：把 $r_d$ 当成电阻画进交流通路。
$r_d$ 与 $I_D$ 成反比，所以换一个静态工作点，交流模型就要重算。}

\subhdr{6.6 三道判断题：该用哪个模型}

（1）\textbf{12 V 汽车灯泡串联二极管防反接}。电源 $12\,\mathrm{V}\gg0.7\,\mathrm{V}$，
用恒压降：$I=(12-0.7)/R$，误差不到 $6\%$；用理想模型也可以，只要不关心灯泡略暗。
（2）\textbf{逻辑门判断 0/1}。只看「高于阈值还是低于阈值」，用理想模型：导通即 $0\,\mathrm{V}$、
截止即开路。此时报具体电压没有意义（§6.4 的电压误差 $100\%$）。
（3）\textbf{精密整流或小信号放大}。信号只有几十 $\mathrm{mV}$，必须用折线模型，
把 $r_d=nU_T/I_D$ 算进去（§6.5）；否则线性化误差会吃掉整个信号。
三题的共同判据：\textbf{先问「题目要的是通断、电流还是电压精度」，再选模型}。

\subhdr{6.7 三种模型都不够的时候}

三种模型都建立在「温度不变、参数已知」的前提上。三种情形必须离开模型、回到指数式或实测：
（1）\textbf{温度跨越大}：$-40$ 到 $+125\,{}^\circ\mathrm{C}$ 时 $U_F$ 变化可达 $0.3\,\mathrm{V}$ 以上（§7.1），
恒压降的 $0.7\,\mathrm{V}$ 不再成立；
（2）\textbf{电流跨数量级}：从 $\mu\mathrm{A}$ 到 A，$U_F$ 从 $0.5$ 变到 $1\,\mathrm{V}$ 以上（§6.2），
单一常数没有意义；
（3）\textbf{要求小信号精度}：必须用 $r_d=nU_T/I_D$（§6.5），而且 $I_D$ 要取实际工作点。
判据一句话：\textbf{模型是「一个数代替一条曲线」，曲线跨得越远，代替得越差}。

\hdr{§7}{温度与材料：用曲线反推参数（p53--p56）}

\subhdr{7.1 两条温度曲线，哪条是 $80\,{}^\circ\mathrm{C}$（p53--p54）}

p53 给两条曲线问 「which one stands for the characteristics of 80\,$^\circ$C and which one is for
20\,$^\circ$C?」，p54 补上标注：红色 $80\,{}^\circ\mathrm{C}$、蓝色 $20\,{}^\circ\mathrm{C}$。
答案是：\textbf{温度高的那条，正向拐点更靠左、反向段下沉更深}。

机制两层：
（1）正向：温度升高 $\Rightarrow n_i$ 指数增大 $\Rightarrow$ 同样电流所需的电压变小。
用 §5.2 的指数式估算（$I_S\propto n_i^2$）：$1\,\mathrm{mA}$ 下，
$20\,{}^\circ\mathrm{C}$ 时 $U\approx0.717\,\mathrm{V}$，$80\,{}^\circ\mathrm{C}$ 时 $U\approx0.543\,\mathrm{V}$，
六十年温差使压降下降约 $0.17\,\mathrm{V}$，平均 $-2.9\,\mathrm{mV}/{}^\circ\mathrm{C}$，
与教材常说的 $-2\sim-3\,\mathrm{mV}/{}^\circ\mathrm{C}$ 一致。
（2）反向：$I_S\propto n_i^2$，$20\to80\,{}^\circ\mathrm{C}$ 使 $n_i^2$ 增大约 $3\times10^{3}$ 倍，
反向段自然下沉得多。

\aside{p18 的「$n\gg p$」与这里的温度效应是同一条链：$n_i$ 一涨，掺杂的优势就被吃掉，
多子浓度、$V_{bi}$、$I_S$ 全都跟着变。}

\subhdr{7.2 硅管还是锗管（p55）}

判别方法已在 §5.5 给出：看拐点电压，不看左右位置。再补一条量级判据：
同样温度下锗管的反向饱和电流比硅管大 $10^3$--$10^5$ 倍（因为 $E_g$ 小、$n_i$ 大）。
所以「反向漏电明显、拐点在 $0.3\,\mathrm{V}$」两条同时成立，就是锗管。

\subhdr{7.3 四个主要参数（p56）}

p56 把四个参数直接标在曲线上：
（1）$I_F$：额定正向平均电流（Rated forward average current）——管子长期允许通过的\emph{平均}电流；
（2）$U_F$：正向压降（Forward voltage drop）——额定电流下的管压降，用来算功耗；
（3）$U_R$：最高反向工作电压（Maximum reverse working voltage）——不是击穿电压，是允许长期承受的反压；
（4）$I_{Rm}$：最大反向电流（Maximum reverse current）——漏电上限。

\par\vspace{2pt}
\noindent\begin{tabular}{@{}>{\raggedright\arraybackslash}p{1.8cm}>{\raggedright\arraybackslash}p{3.4cm}>{\raggedright\arraybackslash}p{3.2cm}>{\raggedright\arraybackslash}p{7.1cm}@{}}
\tabtop
\thead{符号} & \thead{名称} & \thead{关系} & \thead{选型时怎么用}
\\
\tabmid
$I_F$ & 额定正向平均电流 & 由结面积与散热决定 & 实际平均电流 $\le I_F$，通常留 $1.5$--$2$ 倍裕量
\\
$U_F$ & 正向压降 & $0.7\,\mathrm{V}$（Si）或 $0.3\,\mathrm{V}$（Ge） & 算功耗 $P=U_F I_F$，选散热
\\
$U_R$ & 最高反向工作电压 & $U_R\approx(1/2\sim2/3)U_{BR}$ & 反向峰值电压必须小于 $U_R$
\\
$I_{Rm}$ & 最大反向电流 & 与温度强相关 & 判断漏电与发热，高温时要降额
\\
$U_{BR}$ & 反向击穿电压 & 曲线急折处 & 普通管不许工作在击穿区；稳压管专门用这一段
\tabbot
\end{tabular}

\textbf{一个选型算例}（[教材补充：课件只列参数、不给例题]）：某整流电路输出平均电流 $0.5\,\mathrm{A}$，
二极管承受的最高反向电压 $24\,\mathrm{V}$。选 1N4001（$I_F=1\,\mathrm{A}$、$U_R=50\,\mathrm{V}$）：
电流裕量 $1/0.5=2$ 倍，电压裕量 $50/24\approx2.1$ 倍，都合格；
功耗 $P=U_F I_F\approx0.7\times0.5=0.35\,\mathrm{W}$，需要一小片散热。
若换成 1N4007（$1000\,\mathrm{V}$），裕量过大、成本高，属于过设计。

%=============================================================
\textbf{稳压管的三个额外参数}（[教材补充]）：稳定电压 $U_Z$（就是击穿电压）、
动态电阻 $r_Z=\Delta U_Z/\Delta I_Z$（越小稳压越好，几欧到几十欧）、
稳定电流范围 $I_{Z\min}$--$I_{Z\max}$（低于 $I_{Z\min}$ 不稳压、高于 $I_{Z\max}$ 过热）。
选稳压管时先定 $U_Z$，再按负载电流选 $I_{Z\max}$，最后用限流电阻把电流卡在范围内——
这三步是下一册稳压电源电路的入口（本册 p1 的标题 「DC regulated power supply」 就指向那里）。

\subhdr{7.4 热设计：参数表之外的那一半}

选型只对参数还不够，还要算温升。二极管的正向功耗
\[
P=U_F I_F ,
\]
结温
\[
T_j=T_a+P\,R_{\mathrm{th}(j-a)},
\]
其中 $T_a$ 是环境温度、$R_{\mathrm{th}(j-a)}$ 是结到环境的热阻。
仍用 §7.3 的算例：$P=0.35\,\mathrm{W}$、$R_{\mathrm{th}(j-a)}=50\,\mathrm{K/W}$、$T_a=25\,{}^\circ\mathrm{C}$，
则 $T_j=25+17.5=42.5\,{}^\circ\mathrm{C}$，远低于 $125\,{}^\circ\mathrm{C}$ 上限，安全。
若环境升到 $60\,{}^\circ\mathrm{C}$，$T_j=77.5\,{}^\circ\mathrm{C}$，仍安全；若 $P$ 增到 $2\,\mathrm{W}$，
$T_j=125\,{}^\circ\mathrm{C}$ 顶到上限，必须加散热片——这就是「降额」的实际含义：
\textbf{温度越高，允许的 $I_F$ 越小}。p56 的参数是在 $25\,{}^\circ\mathrm{C}$ 下给的，高温下要打折使用。
[教材补充：热阻与降额不在课件中，补它是为了把 p56 的四个参数接到实际使用。]

\subhdr{7.5 温度降额表：把热阻算成允许电流}

由 §7.4 的 $T_j=T_a+U_FI_FR_{\mathrm{th}(j-a)}\le T_{j\max}$ 反解，
\[
I_F\le\frac{T_{j\max}-T_a}{U_F\,R_{\mathrm{th}(j-a)}}.
\]
取 $T_{j\max}=125\,{}^\circ\mathrm{C}$、$U_F=0.7\,\mathrm{V}$、$R_{\mathrm{th}(j-a)}=50\,\mathrm{K/W}$：

\par\vspace{2pt}
\noindent\begin{tabular}{@{}>{\raggedright\arraybackslash}p{3.6cm}>{\raggedright\arraybackslash}p{4.4cm}>{\raggedright\arraybackslash}p{4.4cm}@{}}
\tabtop
\thead{环境温度 $T_a$} & \thead{热极限 $I_F$} & \thead{受额定值限制后}
\\
\tabmid
$25\,{}^\circ\mathrm{C}$ & $2.86\,\mathrm{A}$ & $1.00\,\mathrm{A}$（额定）
\\
$50\,{}^\circ\mathrm{C}$ & $2.14\,\mathrm{A}$ & $1.00\,\mathrm{A}$
\\
$75\,{}^\circ\mathrm{C}$ & $1.43\,\mathrm{A}$ & $1.00\,\mathrm{A}$
\\
$100\,{}^\circ\mathrm{C}$ & $0.71\,\mathrm{A}$ & $0.71\,\mathrm{A}$
\\
$120\,{}^\circ\mathrm{C}$ & $0.14\,\mathrm{A}$ & $0.14\,\mathrm{A}$
\tabbot
\end{tabular}

读表：$75\,{}^\circ\mathrm{C}$ 以下，热极限高于额定值，按 $1\,\mathrm{A}$ 用；$100\,{}^\circ\mathrm{C}$ 以上，
热极限低于额定值，必须降额。这就是数据手册里那条「降额曲线」的来历。

\hdr{§8}{钳位电路：先假设导通，再回头校验}

\subhdr{8.1 把电路认出来（p57）}

p57--p61 是同一张图，改了五个问句：
$+6\,\mathrm{V}$ 经电阻 $R$ 接到公共节点 $U_F$；两只二极管 $D_A$、$D_B$ 的\textbf{阳极都接在
$U_F$ 节点}、\textbf{阴极分别接输入端子 $A$、$B$}；已知 $U_D=0.7\,\mathrm{V}$、$U_A=3\,\mathrm{V}$、$U_B=0\,\mathrm{V}$。
这个接法的意思是：谁的阴极电位低，谁就先把公共节点拉下来。p57 的 (1) 问
「What is the meaning of $U_D$?」——$U_D$ 就是二极管正向导通时两端固定要掉的那 $0.7\,\mathrm{V}$，
它是本题唯一需要记住的器件参数。

p57 左下写 「Not ideal」，页边还写 「Idea or practical?」——后者是 「Ideal or practical?」 的笔误
（应为「理想还是实际」）。这不是内容错误，但读图时容易误解成「想法还是实践」，记一笔。

\subhdr{8.2 四步判据法（课件没给、但每题都要用）}

处理多个二极管同时存在的电路，不能靠直觉，要按下面四步：

\callout{\textbf{四步法}\\
（1）\textbf{假设全截止}：先把所有二极管拿掉，只用电阻求节点电位；\\
（2）\textbf{逐个检查}：算出每只二极管阳极与阴极的电位差 $U_{AK}$，凡 $U_{AK}>U_D$ 的，都「想导通」；\\
（3）\textbf{取最容易的一只}：阴极电位最低的那只先导通（它的 $U_{AK}$ 最大），令它导通并重算节点；\\
（4）\textbf{回头校验}：在导通那只定住的节点电位下，检查其余二极管是否 $U_{AK}<U_D$；\
若仍有一只满足导通条件，改假设它，重复第（3）步。}

\subhdr{8.3 五个追问逐个算（p58--p61）}

\textbf{(2) Which diode will conduct first?（p58）}从 $U_F=0$ 开始往上抬：
$D_B$ 的阴极是 $0\,\mathrm{V}$，$U_F$ 升到 $0.7\,\mathrm{V}$ 时它就满足 $U_{AK}=0.7\,\mathrm{V}$；
$D_A$ 的阴极是 $3\,\mathrm{V}$，$U_F$ 要升到 $3.7\,\mathrm{V}$ 才够。所以 \textbf{$D_B$ 先导通}。

\textbf{(3) When $D_B$ is ON, what will happen to $D_A$?（p59）}$D_B$ 一导通，就把 $U_F$ 钉在
$0+0.7=0.7\,\mathrm{V}$。此时 $D_A$ 的 $U_{AK}=0.7-3=-2.3\,\mathrm{V}<0$，\textbf{$D_A$ 反向截止}。
这就是「争着导通、互相钳制」：先导通的那只把节点电位拿走，另一只就失去了导通条件。

\textbf{(4) How much is the output?（p60）}
\[
U_F=U_B+U_D=0+0.7=0.7\,\mathrm{V}.
\]
校验：$U_A-U_F=3-0.7=2.3\,\mathrm{V}>0$，$D_A$ 反偏 $2.3\,\mathrm{V}$，结论成立。
此时流过 $R$ 的电流 $I_R=(6-0.7)/R$，全部经 $D_B$ 流入 $B$ 端。
若 $R=1\,\mathrm{k\Omega}$，则 $I_R=5.3\,\mathrm{mA}$，$D_B$ 承担全部电流。

\textbf{(5) How much is $U_F$ if the diode is ideal?（p61）}$U_D=0$，同一套推导给出
$U_F=U_B+0=0\,\mathrm{V}$，校验 $U_A-U_F=3\,\mathrm{V}>0$，$D_A$ 仍截止。
两个结果 $0.7\,\mathrm{V}$ 与 $0\,\mathrm{V}$ 的差，就是「实际模型与理想模型」的全部差别。

\subhdr{8.4 为什么不会两只同时导通}

若 $U_A=U_B$，两只确实可能同时导通、平分电流。但只要 $U_A\ne U_B$，假设两只都导通，
则 $U_F$ 同时等于 $U_A+U_D$ 与 $U_B+U_D$，矛盾。物理上，差值 $3\,\mathrm{V}$ 落在指数关系上，
电流比是
\[
\frac{I_A}{I_B}\sim e^{-(U_A-U_B)/(nU_T)}
=e^{-3/0.0336}\approx1.7\times10^{-39},
\]
即高电位那一路的电流小到可以完全忽略——不是「近似截止」，是「数学上就是零」。
这就是为什么判据法可以直接取「最低的那一路」，而不必解联立方程。

\subhdr{8.5 三个易错点}

（1）\textbf{把 $U_F$ 当成 $6\,\mathrm{V}$}：$R$ 上没有电流时才是 $6\,\mathrm{V}$；一旦导通，$U_F$ 由二极管定，不是由电源定。
（2）\textbf{符号弄反}：$U_{AK}$ 是阳极电位减阴极电位，$D_A$ 的阴极是 $A$、不是 $U_F$。
（3）\textbf{忘记 $R$}：$R$ 的作用是限流并把节点往上拉；没有 $R$，$U_F$ 会被电源直接钉死，题就不成立。

%=============================================================
\subhdr{8.6 镜像接法：把「取最低」翻成「取最高」}

把 p57 的两只二极管\textbf{方向全部翻过来}：阳极分别接输入 $A,B$、阴极共接节点 $U_F$，
电阻从 $U_F$ 下拉到地。这就是 p64 的 E2 结构的两输入版本。此时
\[
U_F=\max\left(U_A,U_B\right)-U_D .
\]
为什么：阴极共节点、阳极接输入，输入最高的那一路把节点顶上去；节点一旦被顶到
$U_{\max}-U_D$，其余各路 $U_{AK}=U_i-(U_{\max}-U_D)<U_D$，全部截止。
把它与 §8.3 并排：

\par\vspace{2pt}
\noindent\begin{tabular}{@{}>{\raggedright\arraybackslash}p{3.0cm}>{\raggedright\arraybackslash}p{3.8cm}>{\raggedright\arraybackslash}p{3.2cm}>{\raggedright\arraybackslash}p{4.4cm}@{}}
\tabtop
\thead{接法} & \thead{二极管方向} & \thead{电阻} & \thead{输出}
\\
\tabmid
p57 型 & 阳极共节点、阴极接输入 & 上拉到 $+V$ & $U_{\min}+U_D$（取最低）
\\
p64 型 & 阴极共节点、阳极接输入 & 下拉到地 & $U_{\max}-U_D$（取最高）
\tabbot
\end{tabular}

两种接法的共同点：\textbf{都是让「最不利」的那一路说了算}。数字电路里，
「取最低」对应与逻辑（有一个 0 就输出 0），「取最高」对应或逻辑（有一个 1 就输出 1）。

\subhdr{8.7 钳位电路的电流与功耗}

p60 的结论只给了电压，没给电流。补上：$R$ 上的电流
\[
I_R=\frac{6-0.7}{R}.
\]
若 $R=1\,\mathrm{k\Omega}$：$I_R=5.3\,\mathrm{mA}$，全部流过导通的 $D_B$，
$D_B$ 的功耗 $P_D=0.7\times5.3\,\mathrm{mA}\approx3.7\,\mathrm{mW}$，
$R$ 的功耗 $P_R=I_R^2R\approx28\,\mathrm{mW}$，都很小。
若把 $R$ 改成 $100\,\Omega$：$I_R=53\,\mathrm{mA}$、$P_D\approx37\,\mathrm{mW}$、$P_R\approx0.28\,\mathrm{W}$，
$R$ 需要选 $1/2\,\mathrm{W}$ 以上。\textbf{$R$ 越小，节点被拉得越硬，但功耗越大}——
这是钳位电路设计里唯一的取舍。

\hdr{§9}{二极管逻辑：钳位电路的另一种读法（p62--p69）}

\subhdr{9.1 E1：三路输入取最低（p62--p63）}

p62 的已知：$U_A=0\,\mathrm{V}$、$U_B=4\,\mathrm{V}$、$U_C=6\,\mathrm{V}$，二极管理想，$+12\,\mathrm{V}$ 经 $R$ 上拉到 $U_O$，
三只二极管的阳极共接 $U_O$、阴极分别接 $A,B,C$。按 §8 的判据，阴极电位最低的 $A$ 路先导通，
把 $U_O$ 钉在 $0\,\mathrm{V}$（理想管）。校验另外两路：$U_B-U_O=4\,\mathrm{V}>0$、$U_C-U_O=6\,\mathrm{V}>0$，
都反偏。所以
\[
U_O=\min\left(U_A,U_B,U_C\right)=0\,\mathrm{V}.
\]
p63 的 「(2) How much is the output?」 答案就是 $0\,\mathrm{V}$。

\subhdr{9.2 E2：三路输入取最高（p64）}

p64 把结构翻了过来：三只二极管的\textbf{阳极接输入}、阴极共接 $U_O$，电阻从 $U_O$ \textbf{下拉到地}。
此时输入电位最高的那一路导通，把 $U_O$ 顶上去。已知 $U_A=12\,\mathrm{V}$、$U_B=8\,\mathrm{V}$、$U_C=6\,\mathrm{V}$，
\[
U_O=\max\left(U_A,U_B,U_C\right)=12\,\mathrm{V}.
\]
两页的对比给出一条口诀：\textbf{上拉电阻 + 阴极接输入 $\Rightarrow$ 取最低；下拉电阻 + 阳极接输入 $\Rightarrow$ 取最高}。
前者是与逻辑（AND），后者是或逻辑（OR）。

\subhdr{9.3 E3：把真值表补全（p65--p69）}

p65 给电路与一张只有表头的真值表，p66--p69 每页填一行。四行的完整结果是：

\par\vspace{2pt}
\noindent\begin{tabular}{@{}>{\raggedright\arraybackslash}p{1.5cm}>{\raggedright\arraybackslash}p{1.5cm}>{\raggedright\arraybackslash}p{1.5cm}>{\raggedright\arraybackslash}p{1.5cm}>{\raggedright\arraybackslash}p{2.0cm}>{\raggedright\arraybackslash}p{1.4cm}@{}}
\tabtop
\thead{$U_A$} & \thead{$U_B$} & \thead{$D_A$} & \thead{$D_B$} & \thead{$U_O$} & \thead{逻辑}
\\
\tabmid
$0\,\mathrm{V}$ & $0\,\mathrm{V}$ & ON & ON & $0\,\mathrm{V}$ & 低
\\
$0\,\mathrm{V}$ & $5\,\mathrm{V}$ & ON & OFF & $0\,\mathrm{V}$ & 低
\\
$5\,\mathrm{V}$ & $0\,\mathrm{V}$ & OFF & ON & $0\,\mathrm{V}$ & 低
\\
$5\,\mathrm{V}$ & $5\,\mathrm{V}$ & ON & ON & $5\,\mathrm{V}$ & 高
\tabbot
\end{tabular}

p69 补的正是第四行：两只都 ON、输出 $5\,\mathrm{V}$。四行合起来，$U_O=\min(U_A,U_B)$，
这是一个\textbf{与门}（只有全高才高）。课件 p69 的结论句与这张表一致。

\textbf{非理想时的电平偏移}：若用实际硅管（$U_D=0.7\,\mathrm{V}$），同一电路的输出变成
$U_O=\min(U_A,U_B)+0.7\,\mathrm{V}$：输入 $5/5\,\mathrm{V}$ 时输出 $5.7\,\mathrm{V}$，
输入 $0\,\mathrm{V}$ 时输出 $0.7\,\mathrm{V}$。低电平被抬高 $0.7\,\mathrm{V}$，
这就是二极管逻辑的固有缺点（电平偏移、扇出小、速度慢），也是它后来被 CMOS 取代的原因。
[教材补充：这一句不在课件里，补它是为了回答「既然能实现与门，为什么不用它」。]

\subhdr{9.4 三个电路的关系图}

\par\vspace{2pt}
\noindent\begin{tabular}{@{}>{\raggedright\arraybackslash}p{2.2cm}>{\raggedright\arraybackslash}p{4.4cm}>{\raggedright\arraybackslash}p{4.4cm}>{\raggedright\arraybackslash}p{4.2cm}@{}}
\tabtop
\thead{题号} & \thead{接法} & \thead{输出} & \thead{逻辑}
\\
\tabmid
E1（p62--p63） & 上拉 $+12\,\mathrm{V}$，阴极接输入 & $\min=0\,\mathrm{V}$ & 与
\\
E2（p64） & 下拉到地，阳极接输入 & $\max=12\,\mathrm{V}$ & 或
\\
E3（p65--p69） & 上拉 $+12\,\mathrm{V}$，阴极接输入 & $\min$，真值表见 §9.3 & 与
\tabbot
\end{tabular}

%=============================================================
\subhdr{9.5 公式卡：这一册要用到的全部式子}

\par\vspace{2pt}
\noindent\begin{tabular}{@{}>{\raggedright\arraybackslash}p{3.8cm}>{\raggedright\arraybackslash}p{7.4cm}>{\raggedright\arraybackslash}p{4.2cm}@{}}
\tabtop
\thead{量} & \thead{式子} & \thead{条件/来源}
\\
\tabmid
热电压 & $U_T=kT/q\approx26\,\mathrm{mV}$ & 300 K
\\
本征浓度 & $n_i=\sqrt{N_CN_V}\,e^{-E_g/(2kT)}$ & 热平衡、非简并
\\
质量作用 & $np=n_i^2$ & 热平衡
\\
电导率 & $\sigma=q(n\mu_n+p\mu_p)$ & 漂移为主
\\
内建电势 & $V_{bi}=U_T\ln(N_AN_D/n_i^2)$ & 突变结、全电离
\\
耗尽层宽 & $W=\sqrt{\frac{2\varepsilon_s(V_{bi}-U)}{q}\left(\frac1{N_A}+\frac1{N_D}\right)}$ & 突变结
\\
二极管方程 & $I=I_S\left(e^{U/(nU_T)}-1\right)$ & $|U|\gg nU_T$
\\
动态电阻 & $r_d=nU_T/I_D$ & Q 点附近小信号
\\
折线模型 & $U_D=U_{th}+I_D r_d$ & 正向导通
\\
钳位取最低 & $U_O=U_{\min}+U_D$ & 阳极共节点、上拉
\\
钳位取最高 & $U_O=U_{\max}-U_D$ & 阴极共节点、下拉
\tabbot
\end{tabular}

\subhdr{9.6 易错点：每条都能指回某一页某一步}

\gapitem{（1）把本征当「没有载流子」}{p10 只说 「few in number」，不是零；
$1.5\times10^{10}\,\mathrm{cm^{-3}}$ 仍是每立方厘米一千五百亿个，只是比掺杂少六个数量级（§1.2）。}

\gapitem{（2）以为 $n=p$ 对掺杂硅也成立}{p8 的 $n=p$ 只在\textbf{本征}时成立。
掺杂后 $np=n_i^2$ 仍成立，但 $n\ne p$；p18 的 $n\gg p$ 就是这条的推论（§2.1）。}

\gapitem{（3）把「多子」当成「只有」}{N 型里仍有空穴，只是少；
p42 的反向电流正是少子贡献的，忽略少子就解释不了反向饱和（§4.2、§5.3）。}

\gapitem{（4）把 $V_{bi}$ 当外加电压}{$0.693\,\mathrm{V}$ 是\textbf{内部}电势差，用电压表在两端量不到；
能测到的是它对外加电压的响应（§3.3、§4.1）。}

\gapitem{（5）认为 $0.7\,\mathrm{V}$ 是常数}{p52 写 $0.7\,\mathrm{V}$ 是 mA 级约定；
$0.1\,\mathrm{mA}$ 时约 $0.62\,\mathrm{V}$、$10\,\mathrm{mA}$ 时约 $0.77\,\mathrm{V}$，
温度系数约 $-2.9\,\mathrm{mV}/{}^\circ\mathrm{C}$（§6.2、§7.1）。}

\gapitem{（6）用理想模型报电压}{p51 的理想模型 $U_D=0$，电压误差 $100\%$；
它只能用来判通断，不能用来报电压值（§6.4）。}

\gapitem{（7）钳位时把 $U_F$ 当电源电压}{p57 的 $R$ 上拉，只有所有二极管都截止时 $U_F$ 才等于 $6\,\mathrm{V}$；
一旦导通，$U_F$ 由二极管定（§8.3）。}

\gapitem{（8）钳位时先假设错的那只导通}{p58 要选\textbf{阴极电位最低}的那只先导通；
若先假设 $D_A$，校验会得到 $U_F=3.7\,\mathrm{V}$，而 $D_B$ 的 $U_{AK}=3.7\,\mathrm{V}>0.7\,\mathrm{V}$，
假设自相矛盾（§8.2）。}

\gapitem{（9）逻辑门把 $U_O$ 算成 $12\,\mathrm{V}$}{E1/E3 是上拉、阴极接输入，输出取最低；
只有 E2（下拉、阳极接输入）才取最高（§9.1--§9.3）。}

\gapitem{（10）忘记非理想时的电平偏移}{实际硅管下 $U_O=\min U_i+0.7\,\mathrm{V}$，
低电平被抬到 $0.7\,\mathrm{V}$；理想模型的 $0\,\mathrm{V}$ 只是近似（§9.3）。}

\gapitem{（11）温度升高只想到「更热更导电」}{温度升高 $n_i$ 涨，正向压降\textbf{下降}、反向漏电\textbf{上升}，
拐点左移；两件事方向相反，不能只记一半（§7.1）。}

\gapitem{（12）把击穿当烧毁}{p49 的 $U_{BR}$ 处曲线急折，但只要限流，管子不坏；
稳压管就长期工作在这里（§4.5）。}

\subhdr{9.7 四步用法：拿到一只二极管先做什么}

\callout{\textbf{做题四步}\\
（1）\textbf{判方向}：三角形指向阴极，先看 $U_{PN}=\varphi_P-\varphi_N$ 是否大于 $U_D$；\\
（2）\textbf{选模型}：问通断用理想、问电流用恒压降、问电压精度用折线（§6.4 的误差表）；\\
（3）\textbf{写回路}：把二极管换成所选模型，列回路方程，注意 $U_D$ 的极性；\\
（4）\textbf{回头校验}：算完检查假设——导通的 $I_D>0$、截止的 $U_{AK}<U_D$；不满足就改假设重来。}

\subhdr{9.8 三种门电路的实际输出对比}

把 §9.1--§9.3 的结果按「理想/实际」两栏并排，可以看出电平偏移有多具体：

\par\vspace{2pt}
\noindent\begin{tabular}{@{}>{\raggedright\arraybackslash}p{2.2cm}>{\raggedright\arraybackslash}p{4.6cm}>{\raggedright\arraybackslash}p{3.4cm}>{\raggedright\arraybackslash}p{3.8cm}@{}}
\tabtop
\thead{题号} & \thead{结构} & \thead{理想输出} & \thead{实际输出（$U_D=0.7\,\mathrm{V}$）}
\\
\tabmid
E1 & 上拉 $+12\,\mathrm{V}$、阴极接输入 & $0\,\mathrm{V}$ & $0.7\,\mathrm{V}$
\\
E2 & 下拉到地、阳极接输入 & $12\,\mathrm{V}$ & $11.3\,\mathrm{V}$
\\
E3 & 上拉 $+12\,\mathrm{V}$、阴极接输入 & $\min(U_A,U_B)$ & $\min(U_A,U_B)+0.7\,\mathrm{V}$
\tabbot
\end{tabular}

\textbf{读表}：与门（E1、E3）的低电平被抬高，或门（E2）的高电平被压低——两者都让\textbf{逻辑摆幅}变小。
摆幅在 $5\,\mathrm{V}$ 系统里从 $5\,\mathrm{V}$ 掉到 $4.3\,\mathrm{V}$，噪声容限随之下降。
这就是二极管逻辑「能实现逻辑、但不适合级联」的定量理由。

\subhdr{9.9 级联时电平会跑掉：为什么二极管逻辑只能做一级}

§9.8 说实际输出有 $0.7\,\mathrm{V}$ 偏移。把两级与门级联，偏移会\textbf{累加}：
第 $k$ 级的低电平是 $0.7k\,\mathrm{V}$、高电平是 $5+0.7k\,\mathrm{V}$（以 $5\,\mathrm{V}$ 输入为例）。
（1）低电平：下一级要求输入低于 $V_{IL}\approx0.8\,\mathrm{V}$ 才算 0；
$0.7\times1=0.7\,\mathrm{V}$ 勉强合格，$0.7\times2=1.4\,\mathrm{V}$ 就已经被当成 1——\textbf{两级就失效}；
（2）高电平：要低于上拉电源 $12\,\mathrm{V}$ 才算 1，$5+0.7k<12$ 给出 $k<10$，
不是限制因素。
所以二极管逻辑的真实约束在\textbf{低电平的累加偏移}：它只能级联一级，
这与 p62--p69 只画单级、不画级联是同一个原因。
\aside{这不是「课件没讲」，而是「课件停在单级是合理的」：讲清偏移的累加，
才知道为什么后来的 DTL/TTL 要在二极管后面加三极管做电平恢复。}

\hdr{§10}{缺口清单：课件哪句话跳了步，本讲义补在哪}

以下 21 条按缺口手册的七类框架记录（本册实际出现六类），每条含「类型｜页码｜课件原文短语｜缺了什么｜补法」。

\gapitem{机制缺口}{p7--p8\quad 「The number of holes and free electrons in an intrinsic semiconductor
are equal.」\quad 缺：为什么必然相等（成对产生、成对复合）\quad 补：§1.2 用「断一根键」的微观图像说明。}

\gapitem{机制缺口}{p9\quad 「electrons move towards the higher potential and holes move towards the
lower potential.」\quad 缺：空穴为什么像正电荷、电流方向为何与电子相反\quad 补：§1.3 价电子接力模型。}

\gapitem{机制缺口}{p10\quad 「The carriers of intrinsic semiconductor are few in number.」\quad
缺：「few」是多少、电导率多小\quad 补：§1.2、§1.4 给出 $3.3\times10^{12}$ 分之一与
$\sigma_i\approx4.4\times10^{-6}\,\mathrm{S/cm}$。}

\gapitem{适用区间}{p18\quad 「The number of free electrons $\gg$ The number of holes」\quad
缺：这个不等式的成立条件与失效温度\quad 补：§2.3 给出 $N_D\gg n_i$，并算出约 700 K 失效。}

\gapitem{拼接缺口}{p11\quad 「Intrinsic semiconductor $\to$ Extrinsic semiconductor」\quad
缺：掺杂能把电导率抬高多少\quad 补：§2.4 算出约 $4.9\times10^5$ 倍。}

\gapitem{适用区间}{p28\quad 「Ignore minority carriers」\quad 缺：这句话只在算结区电荷/内建场时成立；
算反向电流时少子恰是主角\quad 补：§3.1 点名，§4.2、§5.3 用少子解释反向饱和电流。}

\gapitem{机制缺口}{p34\quad 「the fixed ions they leave behind establish a 「built-in」 electric field
at the junction」\quad 缺：场的方向、大小与 $V_{bi}$ 的算式\quad 补：§3.3 完整积分推导，
得 $V_{bi}=U_T\ln(N_A N_D/n_i^2)$。}

\gapitem{代数跳步}{p39\quad 「When these two fields add, the resultant field at the junction is smaller
in magnitude」\quad 缺：「相加」加出多少、势垒高度怎么变\quad 补：§4.1 写成 $q(V_{bi}-U)$，
并与 §3.4 的 $W$ 公式对上。}

\gapitem{拼接缺口}{p43 $\to$ p49\quad 「Forward biased\quad ON」 之后直接出现 「Characteristic curve」\quad
缺：从 ON/OFF 突变到平滑曲线的过渡\quad 补：§5.2 用肖克利方程把突变展成指数过渡。}

\gapitem{答案形态}{p50\quad 五个问句 「(1)--(5)」 全页无答案\quad 缺：五个名词的定义与量级\quad
补：§5.2 逐个作答并列成表。}

\gapitem{形式改写}{p52\quad 「Forward bias\quad $U_D=0.7\,\mathrm{V}$(Si)」\quad 缺：$0.7$ 被当成常数，
未说它随电流与温度变\quad 补：§6.2 用指数式给出 $0.1/1/10\,\mathrm{mA}$ 下
$0.619/0.696/0.774\,\mathrm{V}$；§7.1 给温度系数。}

\gapitem{机制缺口}{p53--p54\quad 「which one stands for the characteristics of 80\,$^\circ$C and which
one is for 20\,$^\circ$C?」\quad 缺：温度为什么使拐点左移、反向电流增大\quad 补：§7.1 用
$I_S\propto n_i^2$ 与 $-2.9\,\mathrm{mV}/{}^\circ\mathrm{C}$ 定量说明。}

\gapitem{答案形态}{p55\quad 「which one is for Si-diode and which one is for Ge-diode?」\quad
缺：判别方法与 p52 的 $0.3\,\mathrm{V}$、图上 $0.2\,\mathrm{V}$ 的口径\quad 补：§5.5、§7.2。}

\gapitem{机制缺口}{p57--p59\quad 「Idea or practical?」 与 (1)(2)(3) 只问不答\quad
缺：多二极管电路的判据法\quad 补：§8.2 四步法，§8.3 逐题演算，§8.4 定量说明为何只导一只。}

\gapitem{答案形态}{p62--p64\quad 「Find $U_O$」 未给解\quad 缺：E1/E2 的结果与逻辑功能\quad
补：§9.1、§9.2 给出 $0\,\mathrm{V}$ 与 $12\,\mathrm{V}$，并归纳取最低/取最高口诀。}

\gapitem{拼接缺口}{p45\quad 目录列了 「Photodiode」 与 「Light-emitting diode」\quad
缺：全册只讲普通二极管，光电二极管与 LED 没有下落\quad 补：§1.6 末尾说明它们属于后续册，
本册不展开（不臆造内容）。}

\gapitem{机制缺口}{p48\quad 「By structure: Point contact-type, Junction-type」\quad
缺：这两种结构在结电容、允许电流上的差别\quad 补：§4.7 用结面积与 $C_T$ 的关系说明：
点接触结面积小、电容小、电流小；面结面积大、电容大、电流大。}

\gapitem{适用区间}{p56\quad 四个参数只有名称，没有测试条件\quad
缺：这些参数在什么温度、什么电流下测得\quad 补：§7.4、§7.5 说明参数是 $25\,{}^\circ\mathrm{C}$ 下给的，
高温要降额。}

\gapitem{形式改写}{p62\quad 已知条件写 「Given all the diodes are ideal」，紧接着又问
「(1) Are the diodes ideal or not?」\quad 缺：问句与已知自相矛盾\quad 补：§9.1 按已知条件取理想，
把该问句理解为「请确认本题用的是理想模型」。}

\gapitem{适用区间}{p49\quad 「Barrier potential」 与 「Forward characteristic」 两个标签并排\quad
缺：拐点与导通段的分界没有定义\quad 补：§5.2 用指数式说明「拐点」只是曲率最大处，
没有严格边界，$0.5$--$0.7\,\mathrm{V}$ 都算拐点附近。}

\gapitem{适用区间}{p51\quad 「Ideal characteristic」 在 $U=0$ 处画竖直段\quad
缺：竖直段意味着电流可任意大而压降为零，没说这个假设的适用范围\quad
补：§6.1 说明理想模型只在「电源电压远大于 $0.7\,\mathrm{V}$、只判通断」时成立。}



%=============================================================
\subhdr{缺口类型统计}

\par\vspace{2pt}
\noindent\begin{tabular}{@{}>{\raggedright\arraybackslash}p{3.4cm}>{\raggedright\arraybackslash}p{2.2cm}>{\raggedright\arraybackslash}p{9.4cm}@{}}
\tabtop
\thead{类型} & \thead{条数} & \thead{集中页码}
\\
\tabmid
机制缺口 & 7 & p7、p9、p10、p34、p48、p53、p57
\\
适用区间 & 5 & p18、p28、p49、p51、p56
\\
拼接缺口 & 3 & p11、p43$\to$p49、p45
\\
答案形态 & 3 & p50、p55、p62
\\
形式改写 & 2 & p52、p62
\\
代数跳步 & 1 & p39
\tabbot
\end{tabular}

七类里出现最多的是\textbf{机制缺口}（7 条）：课件习惯「给结论、给图、给问句」，
把中间那一步留给学生自己补。这不是课件写错了，而是它默认你会补——本册最容易卡住的
$V_{bi}$、$0.7\,\mathrm{V}$、反向饱和电流，正好都属于这一类。

\hdr{§11}{综合例题：把前面五条线各做一遍}

五道题都只用到本册的内容，题干是对课件页面的改写（基于本页内容的推断，不是往年题）。
每题给全过程，数字与 §1--§9 的复算一致。

\subhdr{11.1 例一：掺杂硅的三个数（p12--p24）}

\callout{\textbf{题}：一块硅在 300 K 掺入磷，$N_D=5\times10^{15}\,\mathrm{cm^{-3}}$，全部电离。
取 $\mu_n=1350$、$\mu_p=480\,\mathrm{cm^2/(V\cdot s)}$、$n_i=1.5\times10^{10}\,\mathrm{cm^{-3}}$。
求 $n$、$p$、$\sigma$、$\rho$，并问它还是不是 N 型。}

解：N 型，多子为电子，$n\approx N_D=5\times10^{15}\,\mathrm{cm^{-3}}$。由质量作用定律
\[
p=\frac{n_i^2}{n}=\frac{2.25\times10^{20}}{5\times10^{15}}=4.5\times10^{4}\,\mathrm{cm^{-3}}.
\]
空穴项 $p\mu_p$ 比 $n\mu_n$ 小约 $10^{11}$ 倍，忽略，得
\[
\sigma\approx qn\mu_n=1.6\times10^{-19}\times5\times10^{15}\times1350
\approx1.08\,\mathrm{S/cm},
\qquad
\rho=\frac{1}{\sigma}\approx0.93\,\Omega\cdot\mathrm{cm}.
\]
与本征的 $4.4\times10^{-6}\,\mathrm{S/cm}$ 相比，提高约 $2.5\times10^{5}$ 倍。
仍是 N 型：$n/p\approx1.1\times10^{11}\gg1$。
\aside{若温度升到 $600\,\mathrm{K}$，$n_i\approx2.1\times10^{15}$（§1.5），与 $N_D$ 只差 2.4 倍，
$n\approx N_D$ 的近似开始失效；$700\,\mathrm{K}$ 时 $n_i\approx1.3\times10^{16}$，已超过 $N_D$，这块硅退回本征。}

\subhdr{11.2 例二：结的四个参数（p27--p36）}

\callout{\textbf{题}：突变 PN 结，$N_A=10^{17}$、$N_D=10^{15}\,\mathrm{cm^{-3}}$，$A=1\,\mathrm{mm^2}$，
$\varepsilon_s=11.7\varepsilon_0$、$n_i=1.5\times10^{10}\,\mathrm{cm^{-3}}$。
求 $V_{bi}$，以及 $U=0$ 与 $U=-5\,\mathrm{V}$ 时的 $W$、$x_n$、$x_p$ 与 $C_T$。}

解：（1）$V_{bi}$：
\[
V_{bi}=0.0259\ln\frac{10^{17}\times10^{15}}{2.25\times10^{20}}
=0.0259\ln\left(4.44\times10^{11}\right)\approx0.693\,\mathrm{V}.
\]
（2）$U=0$：
\[
W=\sqrt{\frac{2\times1.035\times10^{-12}\times0.693}{1.6\times10^{-19}}
\left(10^{-17}+10^{-15}\right)}\approx0.95\,\mu\mathrm{m}.
\]
\[
x_n=\frac{10^{17}}{1.01\times10^{17}}W\approx0.94\,\mu\mathrm{m},\qquad
x_p\approx0.009\,\mu\mathrm{m}.
\]
（3）$U=-5\,\mathrm{V}$：$V_{bi}-U=5.693\,\mathrm{V}$ 是原来的 8.2 倍，$W\propto\sqrt{V_{bi}-U}$
变为 $\sqrt{8.2}\approx2.86$ 倍，
\[
W\approx2.73\,\mu\mathrm{m},\qquad
C_T=\frac{\varepsilon_sA}{W}=\frac{1.035\times10^{-12}\times0.01}{2.73\times10^{-4}}
\approx38\,\mathrm{pF}.
\]
$U=0$ 时同法得 $C_T\approx109\,\mathrm{pF}$。
\textbf{结论}：反偏 $5\,\mathrm{V}$ 使结电容降到三分之一——这是 §3.7 的定量版。
\aside{注意 $x_n\gg x_p$：耗尽层几乎全在轻掺杂的 N 侧。若把 $N_A,N_D$ 对调，$x_n$ 与 $x_p$ 互换，
但 $V_{bi}$ 与 $W$ 只由乘积 $N_AN_D$ 决定，数值不变。}

\subhdr{11.3 例三：两只二极管串联，三种模型（p49--p52）}

\callout{\textbf{题}：$3\,\mathrm{V}$ 直流源串 $1\,\mathrm{k\Omega}$，再串两只硅二极管（正向）。
用理想、恒压降（每只 $0.7\,\mathrm{V}$）与精确式（$I_S=10^{-12}\,\mathrm{A}$、$n=1.3$）分别求电流。}

解：（1）理想：每只 $U_D=0$，
\[
I=\frac{3}{1000}=3.00\,\mathrm{mA}.
\]
（2）恒压降：每只 $0.7\,\mathrm{V}$，共 $1.4\,\mathrm{V}$，
\[
I=\frac{3-1.4}{1000}=1.60\,\mathrm{mA}.
\]
（3）精确：设每只压降为 $U$，则 $2U=3-1000I$，代入肖克利方程数值求解，
\[
I=1.5765\,\mathrm{mA},\qquad U=0.7118\,\mathrm{V}.
\]
\textbf{三个结果排列}：$3.00>1.60>1.5765\,\mathrm{mA}$。恒压降只比精确值高 $1.5\%$，
理想模型高 $90\%$。两只串联时误差会累加：每只 $0.7\,\mathrm{V}$ 的假设若偏 $0.05\,\mathrm{V}$，
两只就偏 $0.1\,\mathrm{V}$，在 $3\,\mathrm{V}$ 电源里占 $3\%$；串联只数越多，恒压降越要小心。
\aside{本题也说明「串联二极管的压降是叠加的」：$k$ 只串联约 $k\times0.7\,\mathrm{V}$（硅）。
门电路里用二极管做电平移位，靠的就是这个叠加。}

\subhdr{11.4 例四：钳位电路的变式（p57--p61）}

\callout{\textbf{题}：把 p57 的已知改成 $U_A=2\,\mathrm{V}$、$U_B=5\,\mathrm{V}$、$U_D=0.7\,\mathrm{V}$、
上拉 $+6\,\mathrm{V}$。求 $U_F$，并判断两只二极管的状态。}

解：按 §8.2 四步法。
（1）假设全截止：$U_F=6\,\mathrm{V}$（$R$ 上无电流）。
（2）检查：$D_A$ 的 $U_{AK}=6-2=4\,\mathrm{V}>0.7\,\mathrm{V}$；$D_B$ 的 $U_{AK}=6-5=1\,\mathrm{V}>0.7\,\mathrm{V}$，
两只都想导通。
（3）阴极更低的 $D_A$（阴极 $2\,\mathrm{V}$）先导通，把节点钉在
\[
U_F=U_A+U_D=2+0.7=2.7\,\mathrm{V}.
\]
（4）校验 $D_B$：$U_{AK}=2.7-5=-2.3\,\mathrm{V}<0$，截止。
\textbf{答}：$U_F=2.7\,\mathrm{V}$，$D_A$ 导通、$D_B$ 截止。
\aside{若两只输入都改成 $1\,\mathrm{V}$，则 $U_F=1.7\,\mathrm{V}$，两只都导通、共同分担 $R$ 送来的电流；
具体如何分配取决于两管特性差异，理想模型下不能唯一确定——这也是多只并联要均流的原因。}

\subhdr{11.5 例五：温度把工作点推到哪里（p53--p54）}

\callout{\textbf{题}：一只硅二极管工作在 $1\,\mathrm{mA}$，$20\,{}^\circ\mathrm{C}$ 时压降为 $0.717\,\mathrm{V}$。
求 $80\,{}^\circ\mathrm{C}$ 时的压降，并与「$-2\,\mathrm{mV}/{}^\circ\mathrm{C}$」的经验值比较。}

解：用 $I_S\propto n_i^2$ 估算。$20\to80\,{}^\circ\mathrm{C}$ 时
\[
\frac{n_i(353)}{n_i(293)}\approx57,\qquad
\frac{I_S(353)}{I_S(293)}\approx3.3\times10^{3}.
\]
由 $U=nU_T\ln(I/I_S)$，
\[
U(353)=1.3\times0.0304\times\ln\frac{10^{-3}}{10^{-12}\times3.3\times10^{3}}
\approx0.543\,\mathrm{V}.
\]
温度系数
\[
\frac{0.543-0.717}{60}\approx-2.9\,\mathrm{mV}/{}^\circ\mathrm{C},
\]
与常说的 $-2\sim-3\,\mathrm{mV}/{}^\circ\mathrm{C}$ 一致。\textbf{结论}：温度升高，正向压降下降；
若电路靠 $0.7\,\mathrm{V}$ 做基准，$80\,{}^\circ\mathrm{C}$ 时会低 $0.17\,\mathrm{V}$，足以让比较器误判。
\aside{反向方向相反：$I_S$ 涨 $3.3\times10^3$ 倍，漏电急剧增大（§4.6）。
「正向压降变小」与「反向漏电变大」是同一件事的两面，都来自 $n_i$ 增大。}

\subhdr{11.6 例六：把 E1 换成非理想（p62--p63）}

\callout{\textbf{题}：p62 的 E1，已知 $U_A=0\,\mathrm{V}$、$U_B=4\,\mathrm{V}$、$U_C=6\,\mathrm{V}$，
但二极管改为硅管 $U_D=0.7\,\mathrm{V}$，$R=1\,\mathrm{k\Omega}$。求 $U_O$ 与各管状态。}

解：仍是「阴极接输入、阳极共节点」的取最低结构。最低输入 $A=0\,\mathrm{V}$，
$D_A$ 导通后把节点钉在
\[
U_O=U_A+U_D=0+0.7=0.7\,\mathrm{V}.
\]
校验：$D_B$ 的 $U_{AK}=0.7-4=-3.3\,\mathrm{V}<0$，截止；$D_C$ 为 $-5.3\,\mathrm{V}$，截止。
电流 $I_R=(12-0.7)/1000=11.3\,\mathrm{mA}$，全部流过 $D_A$；
$D_A$ 功耗 $P=0.7\times11.3\,\mathrm{mA}\approx7.9\,\mathrm{mW}$，$R$ 功耗 $I_R^2R\approx0.13\,\mathrm{W}$。
\textbf{答}：$U_O=0.7\,\mathrm{V}$，$D_A$ 导通、$D_B/D_C$ 截止；与理想的 $0\,\mathrm{V}$ 差一个 $U_D$。
\aside{这就是 §9.8 的定量版：理想与实际的差别只在「输出被抬高一个 $U_D$」，
通断判断完全一致。}

\subhdr{11.7 例七：E2 的非理想输出（p64）}

\callout{\textbf{题}：p64 的 E2，$U_A=12\,\mathrm{V}$、$U_B=8\,\mathrm{V}$、$U_C=6\,\mathrm{V}$，
二极管为硅管 $U_D=0.7\,\mathrm{V}$，$R=1\,\mathrm{k\Omega}$，求 $U_O$ 与各管状态、电流。}

解：E2 是「阳极接输入、阴极共节点、电阻下拉」的取最高结构。最高的 $A=12\,\mathrm{V}$ 经 $D_A$ 导通，
把节点顶到
\[
U_O=U_A-U_D=12-0.7=11.3\,\mathrm{V}.
\]
校验：$D_B$ 的 $U_{AK}=8-11.3=-3.3\,\mathrm{V}<0$，截止；$D_C$ 为 $-5.3\,\mathrm{V}$，截止。
下拉电阻上的电流 $I_R=11.3/1000=11.3\,\mathrm{mA}$，由 $A$ 端经 $D_A$ 流入地。
$D_A$ 功耗 $0.7\times11.3\,\mathrm{mA}\approx7.9\,\mathrm{mW}$。
\textbf{答}：$U_O=11.3\,\mathrm{V}$，$D_A$ 导通、$D_B/D_C$ 截止。
\aside{与理想的 $12\,\mathrm{V}$ 相比，高电平被压低一个 $U_D$；这正是 §9.8 表里「或门高电平被压低」的具体数。}

\subhdr{11.8 例八：反向漏电随温度涨多少（p53--p56）}

\callout{\textbf{题}：某硅二极管在 $25\,{}^\circ\mathrm{C}$ 时反向漏电 $I_R=1\,\mu\mathrm{A}$，
反向电压 $20\,\mathrm{V}$。按理想扩散电流模型估算 $80\,{}^\circ\mathrm{C}$ 的漏电与功耗。}

解：由 §4.6，理想模型给出 $I_S\propto n_i^2$，$25\to80\,{}^\circ\mathrm{C}$ 约 $3.4\times10^3$ 倍：
\[
I_R(80\,{}^\circ\mathrm{C})\approx1\,\mu\mathrm{A}\times3.4\times10^{3}
\approx3.4\,\mathrm{mA},
\qquad
P=20\,\mathrm{V}\times3.4\,\mathrm{mA}\approx68\,\mathrm{mW}.
\]
这个数已不能忽略（普通塑封管的热阻约 $50\,\mathrm{K/W}$，$68\,\mathrm{mW}$ 会再抬高结温 $3.4\,\mathrm{K}$）。
\textbf{但要说明}：实测漏电通常只放大 $10$--$50$ 倍（约 $10$--$50\,\mu\mathrm{A}$），
因为表面漏电、复合电流等机制不按 $n_i^2$ 走。两个数都指向同一结论：
\textbf{高温下必须按实测漏电（或手册曲线）校核，不能只用室温的 $1\,\mu\mathrm{A}$。}
\aside{这就是 §7.5 降额表存在的理由：额定参数是 $25\,{}^\circ\mathrm{C}$ 的值，
高温下有效值由热平衡决定，而不是由手册的第一行决定。}

\subhdr{11.9 自测十八问（答案在右，遮住右栏自测）}

\par\vspace{2pt}
\noindent\begin{tabular}{@{}>{\raggedright\arraybackslash}p{7.6cm}>{\raggedright\arraybackslash}p{5.6cm}>{\raggedright\arraybackslash}p{2.4cm}@{}}
\tabtop
\thead{问题} & \thead{答案} & \thead{出处}
\\
\tabmid
1. 本征半导体中电子与空穴数目关系？ & 相等，均为 $n_i$ & p8、§1.2
\\
2. N 型半导体的多数载流子是什么？ & 自由电子 & p19、§2.1
\\
3. 磷的第五个电子为什么室温下就能电离？ & 束缚能约 $45\,\mathrm{meV}\ll kT$ & §2.1
\\
4. 少子浓度怎么由多子算？ & $p=n_i^2/N_D$ & §2.1
\\
5. 耗尽层主要向哪一侧扩展？ & 轻掺杂一侧 & §3.2
\\
6. 内建电场方向？ & 由 N 区指向 P 区 & p34、§3.3
\\
7. 内建电势的公式？ & $U_T\ln(N_AN_D/n_i^2)$ & §3.3
\\
8. 正偏时耗尽层变厚还是变薄？ & 变薄 & p40、§4.1
\\
9. 反向电流为什么不随电压变？ & 少子供给由温度决定 & §4.4
\\
10. $0.7\,\mathrm{V}$ 是普适常数吗？ & 不是，mA 级、室温的约定 & §6.2
\\
11. 只有几伏电源、只要通断，用哪个模型？ & 理想模型 & §6.1
\\
12. 要算电流且电源几伏到几十伏，用哪个？ & 恒压降模型 & §6.2
\\
13. 小信号或精度要求高时用哪个？ & 折线模型（含 $r_d$） & §6.3
\\
14. 温度升高，正向压降怎么变？ & 下降，约 $-2\sim-3\,\mathrm{mV}/{}^\circ\mathrm{C}$ & §7.1
\\
15. 钳位电路哪只二极管先导通？ & 阴极电位最低的 & §8.2
\\
16. E1 的理想输出？ & $0\,\mathrm{V}$（取最低） & §9.1
\\
17. E2 的理想输出？ & $12\,\mathrm{V}$（取最高） & §9.2
\\
18. 反向击穿一定会烧管吗？ & 限流则不烧，稳压管就工作在此 & §4.5
\tabbot
\end{tabular}

\hdr{§12}{三句话总结}

\begin{enumerate}
\item \textbf{半导体的一切数值都从 $n_i$ 长出来}：本征浓度 $n_i\propto e^{-E_g/(2kT)}$，
它决定载流子多少、决定 $V_{bi}=U_T\ln(N_A N_D/n_i^2)$，也就决定了硅 $0.7\,\mathrm{V}$、锗 $0.3\,\mathrm{V}$；
\item \textbf{PN 结的「单向」是两个力的平衡被外电压打破}：扩散想过去、内建电场推回来，
正偏压势垒则导通、反偏抬势垒则截止，$0.7\,\mathrm{V}$ 是势垒被压到接近零时的工程约定；
\item \textbf{二极管在电路里有三种替身}：理想开关管通断、恒压降管电流、折线模型管精度，
选哪一种取决于题目问什么——问通断用理想，问电流用恒压降，问电压精度用折线。
\end{enumerate}

%=============================================================
\hdr{§13}{与课件的页码对照表}

\par\vspace{2pt}
\noindent\begin{tabular}{@{}>{\raggedright\arraybackslash}p{1.9cm}>{\raggedright\arraybackslash}p{6.4cm}>{\raggedright\arraybackslash}p{7.5cm}@{}}
\tabtop
\thead{课件页} & \thead{课件内容} & \thead{本讲义位置}
\\
\tabmid
p1 & 章节封面 & §0
\\
p2 & 导体/半导体/绝缘体 & §1.1（补 $E_g$ 与电导率量级）
\\
p3--p5 & 本征定义、8 电子、常温稳定 & §1.2
\\
p6--p8 & 电子—空穴对、$n=p$ & §1.2（补质量作用定律）
\\
p9 & 载流子漂移方向 & §1.3
\\
p10--p11 & 本征导电弱、掺杂出口 & §1.4（补 $\sigma_i$ 计算）
\\
p12--p13 & 掺杂定义、N/P 分类 & §2.1
\\
p14--p19 & 磷 +5、多余电子、多子/少子 & §2.1、§2.3
\\
p20--p24 & 硼 +3、空穴、导电能力 & §2.2、§2.4
\\
p25--p26 & 晶格图判别题 & §2.5
\\
p27--p32 & 两块材料、拼合、扩散与漂移 & §3.1
\\
p33--p34 & 扩散、固定离子 & §3.2
\\
p35--p36 & 漂移、两种效应并存 & §3.3、§3.5
\\
p37--p40 & 正偏、$0.7\,\mathrm{V}$ 导通 & §4.1
\\
p41--p42 & 反偏、耗尽层变厚 & §4.2
\\
p43 & ON/OFF 总结 & §4.3
\\
p44--p48 & 结构、符号、封装、分类 & §5.1
\\
p49--p50 & 伏安曲线、五个问句 & §5.2--§5.4
\\
p51 & 理想模型 & §6.1
\\
p52 & 恒压降模型 & §6.2
\\
p53--p54 & 温度两条曲线 & §7.1
\\
p55 & Si/Ge 判别提问 & §5.5、§7.2
\\
p56 & 四个主要参数 & §7.3
\\
p57--p61 & 钳位电路五问 & §8
\\
p62--p64 & E1、E2 & §9.1、§9.2
\\
p65--p69 & E3 真值表 & §9.3
\tabbot
\end{tabular}

\hdr{附录 A}{数量级速查}

\par\vspace{2pt}
\noindent\begin{tabular}{@{}>{\raggedright\arraybackslash}p{4.4cm}>{\raggedright\arraybackslash}p{4.8cm}>{\raggedright\arraybackslash}p{5.2cm}@{}}
\tabtop
\thead{量} & \thead{数量级} & \thead{备注}
\\
\tabmid
$U_T$（300 K） & $26\,\mathrm{mV}$ & $=kT/q$
\\
$n_i$（Si/Ge） & $10^{10}$ / $10^{13}\,\mathrm{cm^{-3}}$ & 禁带越窄越大
\\
掺杂 $N_D,N_A$ & $10^{14}$--$10^{20}\,\mathrm{cm^{-3}}$ & 跨六个数量级
\\
$V_{bi}$（Si/Ge） & $0.7$ / $0.3\,\mathrm{V}$ & 中等掺杂
\\
$W$（$10^{16}$ 掺杂） & $0.4\,\mu\mathrm{m}$ & 反偏 $5\,\mathrm{V}$ 时约 $1.2\,\mu\mathrm{m}$
\\
$C_T$（$1\,\mathrm{mm^2}$） & $10$--$10^{2}\,\mathrm{pF}$ & 随反偏减小
\\
$I_S$（Si/Ge） & $\mathrm{nA}$ / $\mu\mathrm{A}$ & 强温度依赖
\\
$r_d$（$1\,\mathrm{mA}$） & $30\,\Omega$ & $=nU_T/I_D$
\\
$U_F$（mA 级） & $0.7$ / $0.3\,\mathrm{V}$ & 随电流、温度变
\tabbot
\end{tabular}

\hdr{附录 B}{单位换算与常用记号}

\par\vspace{2pt}
\noindent\begin{tabular}{@{}>{\raggedright\arraybackslash}p{4.4cm}>{\raggedright\arraybackslash}p{6.2cm}>{\raggedright\arraybackslash}p{3.8cm}@{}}
\tabtop
\thead{换算} & \thead{关系} & \thead{用在哪}
\\
\tabmid
长度 & $1\,\mu\mathrm{m}=10^{-4}\,\mathrm{cm}=10^{-6}\,\mathrm{m}$ & $W$、$x_n$、$x_p$
\\
浓度 & $1\,\mathrm{cm^{-3}}=10^{6}\,\mathrm{m^{-3}}$ & $n_i$、$N_D$、$N_A$
\\
迁移率 & $1\,\mathrm{cm^2/(V\cdot s)}=10^{-4}\,\mathrm{m^2/(V\cdot s)}$ & $\mu_n$、$\mu_p$
\\
电导率 & $1\,\mathrm{S/cm}=100\,\mathrm{S/m}$ & $\sigma$
\\
电阻率 & $1\,\Omega\cdot\mathrm{cm}=0.01\,\Omega\cdot\mathrm{m}$ & $\rho$
\\
电容 & $1\,\mathrm{pF}=10^{-12}\,\mathrm{F}$ & $C_T$、$C_D$
\\
热电压 & $U_T=25.9\,\mathrm{mV}$（300 K） & 全部指数式
\\
\tabbot
\end{tabular}

\hdr{附录 C}{常用二极管速查}

以下型号只作「参数长什么样」的例子，具体数值以手册为准。

\par\vspace{2pt}
\noindent\begin{tabular}{@{}>{\raggedright\arraybackslash}p{2.6cm}>{\raggedright\arraybackslash}p{2.6cm}>{\raggedright\arraybackslash}p{4.4cm}>{\raggedright\arraybackslash}p{4.8cm}@{}}
\tabtop
\thead{型号} & \thead{类型} & \thead{主要额定} & \thead{典型用途}
\\
\tabmid
1N4001 & 整流 & $1\,\mathrm{A}$ / $50\,\mathrm{V}$ & 低频小功率整流
\\
1N4007 & 整流 & $1\,\mathrm{A}$ / $1000\,\mathrm{V}$ & 高压整流
\\
1N4148 & 开关 & 约 $100\,\mathrm{V}$ / 几百 $\mathrm{mA}$ & 小信号开关、逻辑
\\
1N5819 & 肖特基 & $1\,\mathrm{A}$ / $40\,\mathrm{V}$，$U_F\approx0.45\,\mathrm{V}$ & 低压大电流、开关电源
\\
1N5408 & 整流 & $3\,\mathrm{A}$ / $1000\,\mathrm{V}$ & 较大电流整流
\tabbot
\end{tabular}

\textbf{读表三点}：肖特基的 $U_F$ 明显低于硅 PN 结（$0.45$ 对 $0.7\,\mathrm{V}$），
所以低压大电流场合选它；开关管的结电容小、反向恢复快，所以逻辑电路选它；
整流管的 $I_F$ 大但速度慢。三者的取舍正好对应 p48 的三条分类标准。
[教材补充：具体型号与额定值不在课件中，只作量级参照。]

\hdr{附录 D}{术语中英对照}
{\small

课件全英文，做题与考试常用中文，这张表用于把两边对上（按本册出现顺序）。

\par\vspace{2pt}
\noindent\begin{tabular}{@{}>{\raggedright\arraybackslash}p{6.8cm}>{\raggedright\arraybackslash}p{8.6cm}@{}}
\tabtop
\thead{英文（课件用词）} & \thead{中文（本讲义用词）}
\\
\tabmid
semiconductor & 半导体
\\
intrinsic semiconductor & 本征半导体
\\
extrinsic semiconductor & 杂质半导体
\\
doping & 掺杂
\\
donor / acceptor & 施主 / 受主
\\
valence electron & 价电子
\\
covalent bond & 共价键
\\
free electron / hole & 自由电子 / 空穴
\\
charge carrier & 载流子
\\
majority / minority carrier & 多数 / 少数载流子
\\
drift motion & 漂移运动
\\
diffusion motion & 扩散运动
\\
depletion region & 耗尽层
\\
built-in electric field & 内建电场
\\
barrier potential & 势垒电压（死区电压）
\\
forward / reverse bias & 正向 / 反向偏置
\\
reverse saturation current & 反向饱和电流
\\
breakdown voltage & 击穿电压
\\
voltage-ampere characteristic & 伏安特性
\\
ideal diode & 理想二极管
\\
constant voltage drop & 恒压降（模型）
\\
dynamic resistance & 动态电阻
\\
anode / cathode & 阳极 / 阴极
\\
rectifier diode & 整流二极管
\\
switching diode & 开关二极管
\\
clamping & 钳位
\\
truth table & 真值表
\\
logic gate & 逻辑门
\\
mobility & 迁移率
\\
conductivity & 电导率
\tabbot
\end{tabular}}

\flowerrule

{\small\color{inkgray}
本讲义所有数值已用 Python 独立复算：$U_T=25.85\,\mathrm{mV}$、$n_i$ 与 $V_{bi}$（Si $0.693\,\mathrm{V}$、
Ge $0.312\,\mathrm{V}$）、$\sigma_i=4.4\times10^{-6}\,\mathrm{S/cm}$、$\sigma_n=2.16\,\mathrm{S/cm}$、
$W=0.42\,\mu\mathrm{m}$、$E_{\max}=3.3\times10^{4}\,\mathrm{V/cm}$、恒压降与折线模型的
$9.300/9.267\,\mathrm{mA}$ 与精确解 $9.229\,\mathrm{mA}$、$r_d$ 的四个量级、
钳位的 $U_F=0.7\,\mathrm{V}$ 与 $0\,\mathrm{V}$、E1/E2/E3 的 $0/12/\min$ 结果。
标 [教材补充] 的各处（能带与材料数值、电离能与冻结、$W/E_{\max}/C_T$、击穿两机制、掺杂范围、折线模型、扩散电容、热阻与降额、器件型号、单位换算、术语对照）不是课件原文，
只作理解与计算支撑，不作文档口径。}
\end{document}
